\documentclass[11pt]{article}
\usepackage[letterpaper,margin=1in]{geometry}
\usepackage[T1]{fontenc}
\usepackage{lmodern,amsmath,amssymb,amsthm,booktabs,microtype,xcolor}
\usepackage[colorlinks=true,linkcolor=blue!45!black,urlcolor=blue!45!black,citecolor=blue!45!black]{hyperref}
\newtheorem{theorem}{Theorem}[section]
\newtheorem{lemma}[theorem]{Lemma}
\newtheorem{proposition}[theorem]{Proposition}
\newtheorem{corollary}[theorem]{Corollary}
\newcommand{\cind}{c_{\mathrm{ind}}}
\newcommand{\bind}{b_{\mathrm{ind}}}
\setlength{\emergencystretch}{2em}
\raggedbottom
\hypersetup{pdftitle={Bridge completion and an equivalent branch-vertex formulation of the induced cycles conjecture},
pdfsubject={Exact reduction and sharp examples; the general conjecture remains unresolved}}
\title{Bridge Completion and an Equivalent Branch-Vertex\\
Formulation of the Induced Cycles Conjecture}
\author{}
\date{October 29, 2026}
\begin{document}
\maketitle

\begin{abstract}
For a subcubic graph $H$, let $\bind(H)$ be the maximum number of its
degree-three vertices contained in an induced $2$-regular subgraph.
We construct a simple cubic completion $\mathcal C(H)$ by attaching a
fixed thirteen-vertex graph at every missing incidence. If
$D=3|V(H)|-2|E(H)|$, then
\[
 |V(\mathcal C(H))|=|V(H)|+13D,\qquad
 \cind(\mathcal C(H))=7D+\bind(H).
\]
Consequently, the cubic induced cycles conjecture is equivalent to
$\bind(H)\geq |E(H)|-|V(H)|$ for every subcubic graph $H$.
The reduction preserves the exact surplus over the conjectured bound.
Further exact formulas describe edge replacements and infinite
families attaining the half bound. A matching-polytope argument
supplies a sufficient condition for nonuniform replacements over
bridgeless cubic bases. Any counterexample containing a bridge would
yield a counterexample obtained by joining two copies of one rooted
graph.
A fractional cut criterion certifies arbitrary triangle-chain
replacements over 2,254 bases with checked rational witnesses.
A conditional version also certifies the sharp twenty-six-vertex
double of the attachment. A smallest counterexample to the branch
formulation would be obtained from a simple cubic base by expanding
each chosen edge once, with fewer than half as many chosen edges as
base vertices. Those chosen edges must form a forest and satisfy
an additional sparsity condition on every subset.

We also reformulate good sets as odd spanning factors without
isolated edges. An odd pseudoforest of this kind would establish
the conjectured bound. We prove a first repair step from every
perfect matching, and a complete factor census through order twenty
verifies repairs using at most two cycle exchanges. Nevertheless,
the proposed monotone repair assertion is false. For every $k\geq3$,
we construct a cubic graph of order $28k-26$ and an odd pseudoforest
with one isolated edge such that every repair must pass through at
least $k$ isolated edges. This barrier is attained, and the shortest
repair has $k$ exchanges. Each graph in this family has induced
cycles vertex number exactly half its order. Thus the obstructions
refute the repair strategy while satisfying the original bound.
The general conjecture remains unresolved.
\end{abstract}

\section{Problem and scope}
All graphs are finite and simple. A graph is \emph{subcubic} if its
maximum degree is at most three. A set $S$ is \emph{good} when the
induced graph on $S$ is $2$-regular; the empty set is allowed.
Write $\cind(G)=\max\{|S|:S\text{ is good in }G\}$.
The conjecture of Henning, Joos, L\"owenstein and Sasse is
\begin{equation}\label{eq:conj}
 \cind(G)\geq |V(G)|/2\quad\text{for every cubic graph }G.
\end{equation}
The original source is \cite{HJLS}; recent papers
\cite{KSS,RRT} still leave the general assertion unresolved.

A previous note in this repository proves \eqref{eq:conj} for graphs
with a square factor. Here we study attachments across bridges.
The structural results below have direct proofs; finite computational
results are identified separately.
Targeted literature searches did not establish publication priority;
no claim of priority is made. In particular, introducing $\bind$ is
a convenient reformulation of the current problem, not a claimed
resolution under new terminology.

\section{An attachment with an exact boundary cost}
Construct a graph $F$ from three disjoint copies $Q_0,Q_1,Q_2$ of
$K_{2,2}$. In $Q_i$ label the two sides
$\{\ell_i^0,\ell_i^1\}$ and $\{r_i^0,r_i^1\}$, and add
\[
 r_i^j\ell_{i+1}^j\qquad(i\in\mathbb Z/3\mathbb Z,\ j\in\{0,1\}).
\]
Thus $F$ has twelve vertices and is cubic. Set $e=r_0^0\ell_1^0$.
Let $P$ be obtained by subdividing $e$ once, with new vertex $z$.
The \emph{root} $z$ has degree two, and all other vertices of $P$
have degree three.

\begin{lemma}\label{lem:cap}
The largest good sets in $P$ containing and avoiding $z$ have orders
seven and six, respectively. Moreover, $V(P)$ can be partitioned into
a good seven-set containing $z$ and a good six-set avoiding $z$.
\end{lemma}
\begin{proof}
We first show that both $F$ and $F-e$ have induced cycles vertex
number six. In either graph, a good set meets a square in zero, two,
three, or four vertices. A singleton or an opposite pair cannot occur:
each such selected vertex has no selected neighbour in the square
and at most one neighbour outside it. A two-set must be adjacent.

If an entire square is selected, its matching neighbours in the other
two squares must be unselected. For $F$ this immediately empties both
other squares, since their remaining vertices form independent pairs.
The same conclusion holds for $F-e$ if the selected square is $Q_2$.
If it is $Q_0$ or $Q_1$, one neighbouring square is still forced empty.
The last square has one forced zero on its side facing the full square;
any possible adjacent pair or triple would need a selected matching
neighbour in the empty square. Hence it too is empty. Thus a four-set
contributes only four vertices in total.

Suppose no square is full, but some square contains three selected
vertices. Its two selected vertices on the same side have internal
degree one. Both therefore require selected matching neighbours.
Neither required matching edge can be the deleted edge.
The neighbouring square contains both vertices on one side; it must
therefore also be a three-set, with that side doubled. These two
three-sets face each other. On their outward sides, the only selected
vertex has internal degree two, so it cannot be adjacent to any further
selected vertex. The third square cannot contain an adjacent two-set,
and cannot form another facing pair of three-sets. It is empty.
The total is six. If no square contains three or four selected
vertices, the total is at most $2+2+2=6$.

For each $j=0,1$, the six vertices
$\{\ell_i^j,r_i^j:i=0,1,2\}$ form an induced cycle in $F$.
The cycle with $j=1$ also survives in $F-e$. This proves the asserted
maxima of six.

A good set in $P$ avoiding $z$ is exactly a good set in $F-e$, so its
maximum is six. If it contains $z$, it contains both endpoints of $e$.
Suppressing $z$ turns it into a good set in $F$ smaller by one, giving
an upper bound of seven. Subdivide the $j=0$ six-cycle to attain seven.
The $j=1$ six-cycle is its vertex-set complement in $P$, proving the
last assertion.
\end{proof}

Every bridge of a graph is absent from every cycle. Consequently, no
bridge can occur as an edge of an induced $2$-regular subgraph.
This elementary fact is essential in the following construction.

\section{Exact cubic completion}
For a subcubic graph $H$, put
\[
 V_3(H)=\{v:d_H(v)=3\},\qquad
 \bind(H)=\max_{S\text{ good in }H}|S\cap V_3(H)|.
\]
Degrees in this definition are measured in $H$, not in the selected
subgraph. Let $n=|V(H)|$, $m=|E(H)|$, and
\[
 d(v)=3-d_H(v),\qquad D=\sum_v d(v)=3n-2m.
\]
At each vertex $v$, attach $d(v)$ disjoint copies of $P$ by joining
$v$ to their roots. Denote the resulting graph by $\mathcal C(H)$.
These new edges are bridges. The graph $\mathcal C(H)$ is simple and
cubic, is connected whenever $H$ is connected, and has order $n+13D$.

\begin{theorem}[Exact completion formula]\label{thm:completion}
For every subcubic graph $H$,
\begin{align}
 \cind(\mathcal C(H))&=7D+\bind(H),\label{eq:completion}\\
 \cind(\mathcal C(H))-\frac{|V(\mathcal C(H))|}{2}
   &=\bind(H)-(m-n).\label{eq:surplus}
\end{align}
\end{theorem}
\begin{proof}
Take any good set in the completion and intersect it with $V(H)$,
obtaining $S$. No attachment bridge can lie in the selected subgraph,
so every selected vertex of $H$ has exactly two selected neighbours
within $H$. Thus $S$ is good in $H$.

If $v\in S$, every root attached to $v$ is unselected, and each
corresponding copy of $P$ contributes at most six vertices.
If $v\notin S$, each attached copy contributes at most seven.
Conversely, these bounds are attainable simultaneously, for any
fixed good $S$, by Lemma~\ref{lem:cap}. There are no edges between
different copies of $P$. The optimum extending $S$ is therefore
\begin{align*}
 |S|+6\sum_{v\in S}d(v)+7\sum_{v\notin S}d(v)
 &=7D+\sum_{v\in S}(1-d(v))\\
 &=7D+\sum_{v\in S}(d_H(v)-2)\\
 &=7D+|S\cap V_3(H)|.
\end{align*}
The last equality uses $d_H(v)\geq2$ for every vertex of a good set.
Maximizing proves \eqref{eq:completion}. Subtracting $(n+13D)/2$
and using $D=3n-2m$ proves \eqref{eq:surplus}.
\end{proof}

\begin{corollary}[Equivalent formulation]\label{cor:equivalent}
The cubic conjecture \eqref{eq:conj} holds if and only if every
subcubic graph $H$ satisfies
\begin{equation}\label{eq:branch}
 \bind(H)\geq |E(H)|-|V(H)|.
\end{equation}
\end{corollary}
\begin{proof}
Apply \eqref{eq:conj} to $\mathcal C(H)$ and use
\eqref{eq:surplus} for the forward implication.
For the reverse implication take $H$ cubic: then $\bind(H)=\cind(H)$
and $m-n=n/2$.
\end{proof}

For a subcubic graph of minimum degree at least two, the right side of
\eqref{eq:branch} is $|V_3(H)|/2$. Thus in that case the statement asks
for induced cycles covering at least half the degree-three vertices;
degree-two vertices are counted with weight zero.
For general subcubic $H$, the form $m-n$ also accounts for vertices
of degree zero or one.

The reduction is quantitative. If $H$ violates \eqref{eq:branch} by
an integer $q>0$, its completion violates \eqref{eq:conj} by exactly
$q$, and has at most $40n$ vertices. This supplies a certified route
from a prospective subcubic counterexample to a cubic one.
It does not establish that either kind of counterexample exists.

The deficit also has an exact interpretation in the complement.
This clarifies why a decomposition into induced cycles and a
pseudoforest would be sufficient, but would require an additional
existence argument.

\begin{lemma}[Complement identity]\label{lem:complement}
If $S$ is good in a subcubic graph $H$ and $T=V(H)\setminus S$, then
\[
 |S\cap V_3(H)|-(m-n)=|T|-|E(H[T])|.
\]
Consequently, \eqref{eq:branch} is equivalent to the existence of a
good set whose complement has at most as many edges as vertices.
A sufficient, potentially stronger condition is that the complement
be a pseudoforest: every component has at most one cycle.
\end{lemma}
\begin{proof}
Each selected degree-three vertex has exactly one neighbour outside
$S$, and each selected degree-two vertex has none. Thus the cut
between $S$ and $T$ has $|S\cap V_3(H)|$ edges. Since
$|E(H[S])|=|S|$, counting all edges and vertices gives the identity.
For a pseudoforest, each component has at most as many edges as
vertices, so the right side is nonnegative.
\end{proof}

\section{Reductions for a smallest counterexample}
The branch formulation admits reductions that preserve the direction
of a counterexample. They do not remove the cubic case.

\begin{lemma}[Pruning and suppression]\label{lem:suppression}
If a counterexample to \eqref{eq:branch} exists, one of minimum order
is connected and has minimum degree at least two. Its degree-two
vertices lie in pairwise vertex-disjoint triangles, each containing
two degree-three vertices.
\end{lemma}
\begin{proof}
Both $\bind$ and $m-n$ are additive over components, so at least one
component of a counterexample is a counterexample.

If $v$ has degree zero or one, put $H'=H-v$. Every good set of $H'$
is good in $H$. Moreover, every vertex of degree three in $H'$ has
degree three in $H$, so $\bind(H')\leq\bind(H)$. If $v$ is a leaf,
then $|E(H')|-|V(H')|=m-n$; if $v$ is isolated, this difference is
$m-n+1$. Thus deleting $v$ preserves a violation.

Now suppose $w$ has degree two and its neighbours $u,v$ are
nonadjacent. Suppress $w$: remove $w$ and its two incident edges and
add $uv$. The resulting simple subcubic graph $H'$ has one fewer
vertex and one fewer edge, and its other vertex degrees are unchanged.
Given a good set $S'$ in $H'$, include $w$ precisely when both $u,v$
belong to $S'$. This replaces a selected edge by a selected path, or
leaves an unselected subdivision vertex. In either case it gives a
good set of $H$ with the same number of degree-three vertices.
Hence $\bind(H')\leq\bind(H)$, while $m-n$ is unchanged.
Again a violation descends to a smaller graph.

We next exclude a pendant cycle $C$: all its vertices have degree two
except a single degree-three vertex $u$, whose external neighbour is
$z$. Remove $V(C)$ to obtain $H'$. The quantity $m-n$ decreases by
one. If a good set $S'$ in $H'$ contains $z$, then $z$ has degree
two in $H'$ and three in $H$, so the same set gains one counted
branch vertex when regarded in $H$. If $z\notin S'$, adjoin the
whole cycle $C$, gaining the branch vertex $u$. In both cases
$\bind(H)\geq\bind(H')+1$, so a counterexample again descends.

Every degree-two vertex is therefore in a triangle whose other
vertices both have degree three: a triangle with two degree-two
vertices would be pendant, and a triangle with three would be an
entire component satisfying equality in \eqref{eq:branch}.
Finally, two distinct such triangles cannot share just one vertex,
as that would require degree at least four. If they share an edge,
its two degree-three endpoints and the two degree-two vertices form
an entire diamond component. That component has $m-n=1$ and a good
triangle containing both degree-three vertices, so is not a
counterexample. Thus the triangles are pairwise vertex-disjoint.
\end{proof}

More generally, subdividing edges cannot decrease $\bind$: lift a
good set by selecting an entire replacement path exactly when its
original edge has both endpoints selected. Subdivision also preserves
$m-n$. Consequently, subdivisions of a graph already satisfying
\eqref{eq:branch} cannot give a counterexample. Suppressing a
degree-two vertex in a triangle would introduce a parallel edge;
the simplicity hypothesis prevents applying the preceding reduction
there without an additional argument.

\begin{lemma}[Consecutive degree-two triangles]\label{lem:triangle-pair}
Suppose a simple subcubic graph $H$ contains two vertex-disjoint
triangles $abw$ and $cdz$. Assume $w,z$ have degree two,
$a,b,c,d$ have degree three, and the remaining edges incident with
these six vertices are precisely $ua,bc,dv$, where $u,v$ are distinct
vertices outside the triangles and are nonadjacent.
Delete the six triangle vertices and add $uv$, obtaining $H'$.
Then
\[
 \bind(H)\geq\bind(H')+2,\qquad
 |E(H)|-|V(H)|=|E(H')|-|V(H')|+2.
\]
Consequently, a smallest counterexample to \eqref{eq:branch} has
no such configuration.
\end{lemma}
\begin{proof}
The assumptions ensure that $H'$ is simple and subcubic; all surviving
vertex degrees are unchanged. Let $S'$ be any good set of $H'$.
If $u,v\in S'$, replace its selected edge $uv$ by the induced path
$uabcdv$, adjoining $a,b,c,d$ and leaving $w,z$ unselected.
This gains four selected branch vertices.
If $u\in S'$ and $v\notin S'$, adjoin the triangle $cdz$.
If $u\notin S'$ and $v\in S'$, adjoin $abw$.
If neither endpoint belongs to $S'$, again adjoin $abw$.
In each of these last three cases the added triangle has no selected
neighbour outside itself, and it gains exactly two branch vertices.
These constructions preserve goodness in every case.

Deleting the configuration removes six vertices and nine edges;
adding $uv$ restores one edge. Thus $m-n$ decreases by two.
The two inequalities show that a counterexample would descend to
the smaller graph $H'$. The distinctness and nonadjacency assumptions
are essential to this reduction as stated; the lemma does not dispose
of its loop or parallel-edge alternatives.
\end{proof}

\begin{lemma}[Deleting a good set and its neighbours]\label{lem:closed-neighbourhood}
Let $S$ be good in a simple subcubic graph $H$, let
$T=N_H(S)\setminus S$, and put $H'=H-(S\cup T)$. Then
\begin{align*}
 \bind(H)&\geq\bind(H')+|S\cap V_3(H)|,\\
 (|E(H)|-|V(H)|)-(|E(H')|-|V(H')|)
   &=\sum_{v\in T}(d_H(v)-1)-|E(H[T])|.
\end{align*}
Consequently, if $S$ is nonempty and its branch count is at least
the second expression, a violation of \eqref{eq:branch} descends
to the smaller graph $H'$.
\end{lemma}
\begin{proof}
The union of $S$ with any good set of $H'$ is good in $H$,
because all neighbours of $S$ outside $S$ were removed.
Surviving vertices can only gain branch weight when viewed in $H$.
Since $H[S]$ has $|S|$ edges, the number of deleted edges is
$|S|+\sum_{v\in T}d_H(v)-|E(H[T])|$. Subtract the number
$|S|+|T|$ of deleted vertices to obtain the identity.
\end{proof}

\begin{lemma}[One expanded parallel edge]\label{lem:house}
Suppose $H$ contains vertices $u,v,a,b,w$ with internal edges
$uv,ua,ab,bv,aw,bw$. Assume $w$ has degree two; the other four
vertices have degree three; and their only external edges are
$ux,vy$, where $x,y$ have degree three. This configuration cannot
occur in a minimum-order counterexample to \eqref{eq:branch}.
\end{lemma}
\begin{proof}
First suppose $x\ne y$ and $xy\notin E(H)$. Delete the five
configuration vertices and add $xy$, obtaining a simple subcubic
$H'$. The quantity $|E|-|V|$ decreases by two and all surviving
degrees are unchanged. If a good set of $H'$ contains both $x,y$,
replace its edge $xy$ by the path $xuvy$, gaining two branch
vertices. Otherwise adjoin the triangle $abw$, also gaining two.

If $x\ne y$ and $xy\in E(H)$, delete the five vertices without
adding an edge. Now $|E|-|V|$ decreases by three. If neither
$x$ nor $y$ is selected in a good set of $H'$, adjoin the induced
four-cycle $uabvu$, gaining four branch vertices. Otherwise adjoin
$abw$, gaining two, and recover at least one further branch weight
at a selected member of $\{x,y\}$, whose degree was reduced from
three to two. Each lift therefore gains at least three.

Finally, if $x=y$, delete the configuration and $x$ together.
They have eight internal edges and one external edge, so $|E|-|V|$
decreases by three. Adjoin $uabvu$ to any good set of the remainder,
gaining four branch vertices. All three cases make a counterexample
descend to a smaller graph.
\end{proof}

\begin{theorem}[A simple base for a smallest obstruction]\label{thm:simple-base}
If \eqref{eq:branch} fails, every minimum-order counterexample has
the following form. Start with a connected simple cubic graph $G$
of order $n$, choose a set $A$ of fewer than $n/2$ edges, and replace
each $uv\in A$ by $ua,ab,bv,aw,bw$, using three fresh vertices.
Thus there is at most one triangle on each base edge. The case
$A=\varnothing$ is allowed.
\end{theorem}
\begin{proof}
By Lemma~\ref{lem:suppression}, the graph $H$ is connected and has
minimum degree two. Its degree-two vertices belong to disjoint
triangles, each with two degree-three vertices. Call vertices
outside these triangles base vertices. If there are none, the
degree-three vertices themselves induce a $2$-regular graph and
give $\bind(H)=|V_3(H)|\geq |E(H)|-|V(H)|$, a contradiction.

Every triangle has exactly two external incidences. Following them
partitions the triangles into chains between base vertices.
Connectedness rules out a separate circular chain. Replacing each
chain by an edge gives a connected cubic multigraph $B$. Ordinary
edges between base vertices are chains of length zero. We next
exclude loops and parallel edges in $B$.

A loop has a chain of $k\geq1$ triangles, since $H$ is simple.
Its branch vertices form an induced cycle $C$ of length $2k+1$.
Outside $C$ are its $k$ degree-two triangle tips and one further
neighbour. The expression in Lemma~\ref{lem:closed-neighbourhood}
is at most $k+2$, whereas $C$ has $2k+1\geq k+2$ branch vertices.
The lemma excludes this loop.

For parallel edges, choose two chains between base vertices $u,v$.
If a length-zero chain exists, include it in the pair. This ensures
that the cycle following the two chains has no chord, including
when there is a third parallel chain. Let $K$ be the total number
of triangles on the chosen chains. Simplicity gives $K\geq1$.
The cycle has $2K+2$ branch vertices; its external neighbourhood
consists of $K$ degree-two tips and at most two further vertices.
The deletion expression is at most $K+4$. If $K\geq2$, the cycle
has at least that many branch vertices, contradicting minimality.
If $K=1$, the two chains are one ordinary edge and one triangle,
exactly the configuration of Lemma~\ref{lem:house}. Its external
neighbours have degree three: a base vertex cannot be adjacent to
a degree-two triangle tip. That lemma gives the contradiction.

Hence $B$ is simple. A chain of length at least two contains a pair
to which Lemma~\ref{lem:triangle-pair} applies. If the chain has
length two, the outside endpoints are nonadjacent because $B$ has
no parallel edge; for a longer chain, the next triangle endpoint
has no edge back to the first base vertex. Thus every chain has
length zero or one, as claimed.

Write $k=|A|$. The union of the $k$ gadget triangles is good, since
different triangles are separated by base vertices, and has $2k$
branch vertices. Also $|E(H)|-|V(H)|=n/2+k$. If $k\geq n/2$ this
union already proves the required bound. Therefore $k<n/2$.
\end{proof}

This is a necessary form for a smallest obstruction, not its
exclusion. In particular, the unrestricted simple cubic case
remains when $A=\varnothing$.

\begin{proposition}[Sparsity of the expanded edges]\label{prop:marked-sparsity}
Let $H_A$ be a minimum-order counterexample in the form supplied by
Theorem~\ref{thm:simple-base}. For every nonempty $B\subseteq A$,
put $U=V(B)$ and $q_B=|E(G[U])\setminus A|$. Then
\begin{equation}\label{eq:marked-sparsity}
 2|B|+q_B\leq2|U|-1.
\end{equation}
In particular, $A$ is a forest. Every nontrivial component of this
forest induces in $G$ a tree or a unicyclic graph.
\end{proposition}
\begin{proof}
Let $S$ be the union of the gadget triangles for $B$. These
triangles have no edges between them, so $S$ is good and contains
$2|B|$ branch vertices. Its outside neighbourhood in $H_A$ is
exactly $U$, consisting of base vertices of degree three.
The edges of $H_A[U]$ are precisely $E(G[U])\setminus A$.
Lemma~\ref{lem:closed-neighbourhood} therefore deletes $S\cup U$
at cost $2|U|-q_B$ in $|E|-|V|$, while its good-set lift gains
at least $2|B|$ branch vertices. Minimality forces
$2|B|<2|U|-q_B$, proving \eqref{eq:marked-sparsity}.

If $A$ contained a cycle, take its edge set as $B$. Then
$|B|=|U|$, contradicting the inequality. Hence $A$ is a forest.
For the edge set of any nontrivial component, $|U|=|B|+1$,
so $q_B\leq1$. The induced graph on this component's vertices
is its marked tree together with at most one further edge.
\end{proof}

The entire collection of subset inequalities has an equivalent
description on a smaller auxiliary multigraph. This observation
does not require $H_A$ to be a counterexample.

\begin{lemma}[Contracting the marked forest]\label{lem:marked-quotient}
For a simple cubic graph $G$ and $A\subseteq E(G)$, condition
\eqref{eq:marked-sparsity} for every nonempty $B\subseteq A$
is equivalent to the following two conditions. First, $A$ is a
forest. Second, form a multigraph $J$ whose vertices are the
nontrivial components of that forest, with one edge for each
unmarked edge of $G$ having both endpoints in $V(A)$, between
their components. Retain loops and parallel edges. Then
\[
 |E(J[W])|\leq2|W|-1\qquad
 \text{for every nonempty }W\subseteq V(J).
\]
A loop counts as one edge in this inequality. The empty marking
satisfies both descriptions vacuously.
\end{lemma}
\begin{proof}
The cycle argument above proves necessity of the forest condition
without using minimality. If $W$ is a collection of its components,
take $B$ to contain all their edges. Then
$|V(B)|-|B|=|W|$ and $q_B=|E(J[W])|$, giving necessity of the
displayed inequality.

Conversely, take nonempty $B\subseteq A$, let $c$ be the number
of components of $(V(B),B)$, and let $W$ consist of the components
of $A$ that meet $B$. Since $A$ is a forest, $|V(B)|=|B|+c$
and $c\geq|W|$. Every edge counted by $q_B$ is among the edges
of $J[W]$. Hence
$q_B\leq2|W|-1\leq2c-1$, which is
\eqref{eq:marked-sparsity}.
\end{proof}

The forest restriction does not assert that $A$ is a matching.
Nor do these conditions exclude $A=\varnothing$: they leave the
unrestricted cubic case unresolved.

The following exact calculation records a separate constraint on a
possible obstruction built across a bridge. It helps distinguish
rooted bounds from the unrooted conjecture.

\begin{lemma}[Doubling a rooted attachment]\label{lem:double}
Let $Q$ have one vertex $r$ of degree two and all other vertices of
degree three. Put $q=|V(Q)|$, and let $A$ and $B$ be the maximum
orders of good sets containing and avoiding $r$, respectively.
Use $A=-\infty$ if no good set contains $r$.
Join the roots of two disjoint copies of $Q$ by an edge to form $D(Q)$.
Then $q$ is odd, $D(Q)$ is cubic, and
\[
 \cind(D(Q))=\max\{2B,A+B\}.
\]
In particular, $D(Q)$ is a counterexample if and only if
$2B<q$ and $A+B<q$.
\end{lemma}
\begin{proof}
The degree sum of $Q$ is $3q-1$, so $q$ is odd. The added edge is a
bridge and cannot occur in a selected cycle. At most one of its
endpoints can therefore be selected. The two possible optimum
cardinalities are $2B$ and $A+B$. Each finite candidate value is
attained by the corresponding choices in the two copies.
\end{proof}

For example, if $Q$ is obtained by subdividing an edge $e$ of a cubic
graph $G$, then $B=\cind(G-e)$ and
$A=1+\max\{|S|:S\text{ is good in }G,\ e\subseteq S\}$,
with the same infeasibility convention. Notice that deleting $e$
can create new good sets, so $\cind(G-e)$ cannot simply be replaced
by $\cind(G)$ in this test.

\begin{theorem}[Bridge symmetrization]\label{thm:bridge-symmetry}
Let $G$ be a connected cubic graph with a bridge. Deleting the bridge
gives rooted graphs $Q_1,Q_2$, with roots at its endpoints. Write
$q_i=|V(Q_i)|$, and let $A_i,B_i$ be their root-in and root-out
good-set optima. Define
\[
 \delta_i=\max\{0,A_i-B_i\},\qquad
 d_i=\cind(D(Q_i))=2B_i+\delta_i,
\]
where an infeasible root-in state has $A_i=-\infty$ and
$\delta_i=0$. Then
\begin{equation}\label{eq:bridge-symmetry}
 \cind(G)=\frac{d_1+d_2+|\delta_1-\delta_2|}{2}.
\end{equation}
Consequently, if $G$ violates \eqref{eq:conj}, then at least one
of $D(Q_1),D(Q_2)$ also violates it.
\end{theorem}
\begin{proof}
An induced union of cycles cannot use the bridge, so at most one
root is selected. Each allowable pair of local good sets gives a
global good set. Thus
\[
 \cind(G)=\max\{B_1+B_2,A_1+B_2,B_1+A_2\}
 =B_1+B_2+\max\{\delta_1,\delta_2\}.
\]
Substitute
$2\max\{\delta_1,\delta_2\}=\delta_1+\delta_2+
|\delta_1-\delta_2|$ and apply Lemma~\ref{lem:double}.
If both $d_i\geq q_i$, equation~\eqref{eq:bridge-symmetry} gives
$\cind(G)\geq(q_1+q_2)/2$. Taking the contrapositive proves the
last assertion. Each $D(Q_i)$ is a simple cubic graph of order $2q_i$.
\end{proof}

Thus searching for a counterexample with a bridge can be restricted
to symmetric joins without losing existence. This is not an
order-preserving reduction: the new counterexample can be larger
than $G$. It also places no restriction on a bridgeless counterexample.
In particular, it does not itself prove the original conjecture.

\begin{proposition}[One edge replaced by a triangle]\label{prop:one-triangle}
Let $G$ be cubic of order $n=2h$, and let $e=uv$. Replace $uv$ by
the edges $ua,ab,bv,aw,bw$, where $a,b,w$ are new vertices, obtaining
a subcubic graph $H$. Define
\[
 B_e=\cind(G-e),\qquad
 T_e=\max\{|S|:S\text{ good in }G,\ u,v\in S\},
\]
and $C_e=\max\{|S|:S\text{ good in }G,\ u,v\notin S\}$.
As usual an infeasible maximum is $-\infty$. Then
\begin{equation}\label{eq:one-triangle}
 \bind(H)=\max\{B_e,T_e+2,C_e+2\},\qquad
 |E(H)|-|V(H)|=h+1.
\end{equation}
If $\cind(G)\geq h$ but $H$ violates the branch inequality, then
\[
 \cind(G)=h,\qquad B_e\leq h,\qquad T_e,C_e\leq h-2.
\]
In particular every maximum good set of $G$ contains exactly one
endpoint of $e$. Subdividing $e$ once to obtain $Q$ then gives a
cubic counterexample $D(Q)$ of order $2n+2$ with
$\cind(D(Q))\leq n$.
\end{proposition}
\begin{proof}
Consider a good set of $H$. If it contains $w$, it must contain
$a,b$, and it must avoid $u,v$. The triangle contributes two branch
vertices and the remaining vertices form a good set of $G$ avoiding
both endpoints. If $w$ is unselected, either both $a,b$ are
unselected, or both are selected. The former case is exactly a good
set of $G-e$. In the latter case $u,v$ must be selected, and
contracting the path $uabv$ to $uv$ gives a good set of $G$ containing
both endpoints. Conversely, each of these three kinds of good set
extends as described. This proves the maximum formula. The operation
adds three vertices and four edges, proving the second identity.

If the branch inequality fails, integrality in
\eqref{eq:one-triangle} gives $B_e\leq h$ and $T_e,C_e\leq h-2$.
Every good set of $G$ either contains both endpoints and has size at
most $T_e$, or remains good after deleting $e$ and has size at most
$B_e$. Thus $\cind(G)\leq h$; the hypothesis gives equality.
A maximum good set cannot contain both endpoints or neither endpoint,
by the bounds on $T_e,C_e$.

For the subdivided graph $Q$, the two optima in
Lemma~\ref{lem:double} are $A=T_e+1\leq h-1$ and $B=B_e\leq h$.
Thus $A+B\leq n-1$ and $2B\leq n$. Lemma~\ref{lem:double} gives
$\cind(D(Q))\leq n$, below its half bound $n+1$.
\end{proof}

\begin{proposition}[Simultaneous replacements]\label{prop:simultaneous}
Let $G$ be a simple cubic graph of order $n$, let $A\subseteq E(G)$,
and perform the edge-triangle replacement of
Proposition~\ref{prop:one-triangle} on each edge of $A$, using disjoint
new vertices. Denote the resulting subcubic graph by $H_A$.
For an Eulerian edge set $F\subseteq E(G)$, let $R=V(F)$ be its
support, put $X=E(G[R])\setminus F$, and let $\delta(R)$ be the
edge boundary of $R$. Include $F=\varnothing$. Then
\begin{equation}\label{eq:simultaneous}
 \bind(H_A)=
 \max_{\substack{F\text{ Eulerian}\\X\subseteq A}}
 \bigl\{|R|+2\bigl(|A|-|A\cap\delta(R)|-|X|\bigr)\bigr\}.
\end{equation}
In particular, a connected base requires only $2^{n/2+1}$ Eulerian
states, independently of the number of expanded edges.
\end{proposition}
\begin{proof}
In a cubic graph an Eulerian edge set has degree zero or two at
every vertex. Thus $X$ is a matching: each vertex of its support
already has two incident edges in $F$.

Project a good set of $H_A$ onto the original vertices, obtaining
$R$. An expanded edge has precisely the three local possibilities
described in Proposition~\ref{prop:one-triangle}: an empty gadget,
the path through its two branch vertices, or its whole triangle.
Let $F$ consist of the unexpanded edges with both endpoints selected
and the expanded edges whose path is selected. Every selected
original vertex has degree two in $F$, so $F$ is Eulerian and has
support $R$. Every edge of $X$ must have been expanded and must use
an empty gadget. A marked edge crossing $\delta(R)$ likewise has
an empty gadget. The remaining marked edges can contribute two
branch vertices apiece, giving the upper bound in
\eqref{eq:simultaneous}.

Conversely, take any Eulerian $F$ with $X\subseteq A$. Select $R$;
on marked edges of $F$ select the two branch vertices, and on marked
edges outside $R$ select the entire triangle. Leave the other
gadgets empty. Every selected degree is two, and the branch count
is exactly the displayed expression. This proves equality. The
state count is the dimension formula for the binary cycle space.
\end{proof}

\begin{theorem}[A fractional cut certificate]\label{thm:cut-certificate}
Let $G$ be a simple cubic graph of order $n$. Suppose a probability
distribution on good sets has edge-cut probabilities
$p_e=\Pr(e\in\delta(S))$ satisfying
\begin{equation}\label{eq:cut-certificate}
 \sum_{e\in E(G)}\min\{p_e,1-p_e\}\geq n/2.
\end{equation}
Then $\bind(H_A)\geq n/2+|A|$ for every $A\subseteq E(G)$.
Consequently every cubic completion $\mathcal C(H_A)$ satisfies
the original conjecture. More quantitatively, if
$\tau=\sum_e\min\{p_e,1-p_e\}$, then
$\bind(H_A)\geq |A|+\lceil\tau\rceil$.
\end{theorem}
\begin{proof}
Every good set $S$ of $G$ lifts to a good set of $H_A$ with
branch count $|S|+2|A\setminus\delta(S)|$. Since $G$ is cubic,
each selected vertex has exactly one neighbour outside $S$, so
$|\delta(S)|=|S|$. The expected surplus of this lift over
$n/2+|A|$ is therefore
\begin{align*}
 \mathbb E|S|-n/2+|A|-2\sum_{e\in A}p_e
 &= -n/2+\sum_{e\notin A}p_e+\sum_{e\in A}(1-p_e)\\
 &\geq -n/2+\sum_e\min\{p_e,1-p_e\}\geq0.
\end{align*}
Some lift has nonnegative surplus. Also
$|E(H_A)|-|V(H_A)|=n/2+|A|$, so Theorem~\ref{thm:completion}
gives the assertion for cubic completions.
The same calculation bounds the expected branch count below by
$|A|+\tau$; some lift attains at least this expectation, and
rounding up gives the quantitative assertion. In particular,
$\tau>n/2$ gives a strict unit surplus. Reversing pairs of triangle
replacements preserves this surplus for arbitrary chain lengths.
\end{proof}

The hypothesis is sufficient, not necessary. For the sharp
twenty-six-vertex graph formed from two copies of $P$ joined at
their roots, it cannot hold. If $p$ is the probability that the
joining bridge crosses the selected cut, Lemma~\ref{lem:cap}
gives $\mathbb E|S|\leq12+p$: sets selecting neither root have
size at most twelve, and those selecting one have size at most
thirteen. Therefore every distribution satisfies
\[
 \sum_e\min\{p_e,1-p_e\}
 \leq12+p-\max\{0,2p-1\}\leq25/2<13.
\]
The graph itself satisfies the original conjecture with equality.
This explains why the certificate cannot settle the unrestricted
problem by itself.

\begin{theorem}[Conditional cut certificates]\label{thm:conditional-cut}
Let $G$ be simple cubic of order $n$, and let $P_0,Q_0\subseteq E(G)$
be disjoint sets of forced marked and forced unmarked edges.
Consider a probability distribution on Eulerian edge sets $F$ such
that every chord set $X=E(G[V(F)])\setminus F$ is contained in $P_0$.
Write $p_e=\Pr(e\in\delta(V(F)))$. If
\begin{equation}\label{eq:conditional-cut}
 \sum_{e\in P_0}(1-p_e)+\sum_{e\in Q_0}p_e+
 \sum_{e\notin P_0\cup Q_0}\min\{p_e,1-p_e\}\geq n/2,
\end{equation}
then every marking $P_0\subseteq A\subseteq E(G)\setminus Q_0$
satisfies $\bind(H_A)\geq n/2+|A|$.
\end{theorem}
\begin{proof}
Put $R=V(F)$ and $C=\delta(R)$. Counting the unused incidence at
each vertex of $R$ gives $|R|=2|X|+|C|$. Thus the surplus of the
lift in Proposition~\ref{prop:simultaneous} is exactly
\begin{equation}\label{eq:cut-distance}
 |R|+2(|A|-|A\cap C|-|X|)-(n/2+|A|)
   =|A\mathbin{\triangle}C|-n/2.
\end{equation}
Every state in the distribution is feasible because $X\subseteq
P_0\subseteq A$. The expected symmetric-difference size is
$\sum_{e\in A}(1-p_e)+\sum_{e\notin A}p_e$, bounded below by
the left side of \eqref{eq:conditional-cut}. Hence some feasible
lift has the desired surplus.
\end{proof}

The earlier criterion is the case $P_0=Q_0=\varnothing$.
Unlike it, the conditional version may use Eulerian supports with
chords, provided those chords are forced to be expanded. For the
sharp double of $P$, split on the joining bridge $e$: use
$(P_0,Q_0)=(\varnothing,\{e\})$ or $(\{e\},\varnothing)$.
Exact rational distributions with ten states each attain thirteen
in both cases. These two certificates cover every subset of its
thirty-nine edges. They are checked without an optimization package
in Section~\ref{sec:round16}; a solver status is not used as proof.

\begin{corollary}[Arbitrary triangle chains]\label{cor:triangle-chains}
Suppose $G$ has a distribution satisfying
\eqref{eq:cut-certificate}, or has conditional certificates whose
covered markings together include every $A\subseteq E(G)$.
Replace each edge $uv$ by any number
$k_e\geq0$ of triangles in series: for $k_e>0$, use
triangles $a_i b_i w_i$ with connecting edges
$ua_1,b_1a_2,\ldots,b_{k_e-1}a_{k_e},b_{k_e}v$.
Use distinct new vertices for different chains.
Let $K=\sum_e k_e$ and let $H$ be the resulting subcubic graph.
Then
\[
 \bind(H)\geq n/2+K,\qquad
 |V(\mathcal C(H))|=n+16K,
\]
and $\mathcal C(H)$ satisfies the original half bound.
\end{corollary}
\begin{proof}
Repeatedly remove two consecutive triangles using
Lemma~\ref{lem:triangle-pair}. The outside endpoints are distinct
and nonadjacent because the base is simple and each replacement
is a chain with disjoint internal vertices. This leaves one
triangle on the edges with odd $k_e$ and none on the others.
Theorem~\ref{thm:cut-certificate} or
Theorem~\ref{thm:conditional-cut} applies to this reduced graph.
Reversing each reduction restores two units of branch count and
two units of $|E|-|V|$. Hence $\bind(H)\geq n/2+K$.
The $K$ vertices $w_i$ are the only degree-two vertices, so the
completion formula gives order $n+3K+13K=n+16K$ and the half bound.
\end{proof}

The fixed attachment also admits certificates retaining its root
state. Use labels $\ell_i^j=4i+j$, $r_i^j=4i+2+j$, and root $12$.
The following two lists each contain four good sets with six branch
vertices:
\[
\begin{array}{c|l}
\text{root avoided}&
\{4,6,7,8,9,10\},\ \{5,6,7,8,9,11\},\\
&\{0,1,2,8,10,11\},\ \{0,1,2,9,10,11\}\\[2pt]
\text{root selected}&
\{1,2,3,4,5,6,12\},\ \{0,2,3,4,5,7,12\},\\
&\{0,2,4,6,8,10,12\},\ \{1,2,4,7,9,11,12\}.
\end{array}
\]
In either uniform distribution every edge crosses at most two of
the four selected cuts. These claims follow directly by checking
the nineteen edges of $P$.

\begin{proposition}[Preserving both attachment states]\label{prop:rooted-expansions}
Replace the edges of $P$ by triangle chains with total length $K$.
The resulting subcubic graph has, separately with its root selected
and avoided, a good set containing at least $6+K$ branch vertices.
Attach a fresh copy of $P$ at each of the $K$ new degree-two tips,
leaving the original root of degree two. The resulting rooted graph
$Q$ has order $13+16K$ and has root-in and root-out optima satisfying
\[
 A(Q)\geq7+8K,\qquad B(Q)\geq6+8K.
\]
\end{proposition}
\begin{proof}
First expand a subset $A$ of edges just once. For either of the
displayed four-set distributions, the ordinary good-set lift has
expected branch count
$6+2\sum_{e\in A}(1-p_e)\geq6+|A|$, since $p_e\leq1/2$.
Its original root state is unchanged. Remove pairs of consecutive
triangles to handle arbitrary chains, and reverse the reductions
as in Corollary~\ref{cor:triangle-chains}. The lifts retain the root
state and give branch count at least $6+K$.

For any such good set, completing a new tip contributes seven from
its attached cap if the tip is absent, or six plus the selected tip
if it is present. The combined contribution is seven in either
case. Thus the completed set has size equal to its branch count,
plus $7K$, plus the indicator that the original root is selected.
The asserted lower bounds follow. They are not claims of equality
for the expanded attachment.
\end{proof}

\section{Sharp graphs with bridges}
Let $T$ be a tree in which each vertex has degree one or three,
and let $t$ be its number of degree-three vertices. Replace each such
vertex by a triangle, using its three vertices as distinct attachment
points for the three incident edges of $T$. Replace each leaf by a
copy of $P$, using its root as the attachment point. Retain the edges
corresponding to $T$ as bridges. Write $G_T$ for the resulting graph.
There are $t+2$ leaves, so
\[
 |V(G_T)|=3t+13(t+2)=16t+26.
\]

\begin{theorem}[Exact equality family]\label{thm:trees}
For every such tree,
\[
 \cind(G_T)=8t+13=\frac{|V(G_T)|}{2}.
\]
In addition, $V(G_T)$ partitions into two good sets of this size.
\end{theorem}
\begin{proof}
For $t=0$, the tree is an edge joining two leaves. The graph consists
of two copies of $P$ joined at their roots. At most one root is
selected, so the optimum is $7+6=13$. The complementary sets from
Lemma~\ref{lem:cap}, with opposite choices in the two copies,
give the required partition.

Now let $t\geq1$. The subgraph $T_0$ induced by the nonleaves of $T$
is a tree on $t$ vertices. Because attachment edges are bridges,
a good set selects either an entire internal triangle or none of it.
The selected triangles form an independent set $X$ of $T_0$.
If $v$ has $\ell_v$ leaf neighbours in $T$, selecting its triangle
gains three vertices and reduces the available contribution from
each attached leaf copy from seven to six.
For a fixed $X$, the optimum is therefore
\[
 7(t+2)+\sum_{v\in X}(3-\ell_v)
   =7(t+2)+\sum_{v\in X}d_{T_0}(v).
\]
For any independent $X$, the last sum counts each edge of $T_0$
at most once, so it is at most $t-1$.
Either bipartition class attains $t-1$. This proves the maximum
$7(t+2)+t-1=8t+13$.

For the partition statement, assign the triangles according to a
bipartition of $T_0$. At every leaf copy, assign its root-containing
seven-set to the class opposite its neighbouring triangle, and its
complementary six-set to the triangle's class.
The two resulting sets partition all vertices and are both good:
every attachment edge joins opposite classes. The preceding degree
count shows that both sets have $8t+13$ vertices.
\end{proof}

These orders are $2$ modulo $4$, so these graphs have no square factor.
The theorem thus provides sharp examples outside the class treated
in the preceding note. The partition property is a familiar sufficient
condition for half-coverage: in a cubic graph, a partition into two
induced $2$-regular subgraphs has a perfect matching between its parts.
The equality upper bound here follows from the attachment calculation,
not merely from existence of that partition.

\begin{corollary}[Decorated bridge trees]\label{cor:decorated-trees}
In any graph $G_T$, replace any of its bridges and any edges inside
its leaf copies of $P$ by arbitrary triangle chains. Let $K$ be the
total number of new triangles, and attach a fresh copy of $P$ at
each new degree-two tip. The resulting cubic graph has order
$16t+26+16K$ and induced cycles vertex number at least
$8t+13+8K$.
\end{corollary}
\begin{proof}
Expanding a bridge by one triangle and completing its tip subdivides
the corresponding tree edge by a new degree-three tree vertex,
with a new leaf attached there. Repeating this handles all bridge
chains and produces another tree of the same type. An expanded
original leaf cap with $K_i$ new triangles is the rooted graph of
Proposition~\ref{prop:rooted-expansions}. In either root state it
contributes at least $8K_i$ more vertices than the original cap.
Use a half-size good set from Theorem~\ref{thm:trees} on the tree
after bridge expansions, and replace the states in its expanded
leaf caps by these witnesses. Root states, and hence validity across
bridges, are preserved. Each operation adds sixteen vertices and
at least eight selected vertices, giving the assertion. Internal
edges of the original degree-three-vertex triangles are not among
the edges allowed by this corollary.
\end{proof}

\section{An exact edge replacement}
This section permits loops and parallel edges in an auxiliary base
multigraph $B$; a loop contributes two to its vertex degree. The
constructed graphs remain simple. Let $\nu_2(B)$ be the maximum
number of vertices covered by an ordinary $2$-regular subgraph of
$B$. Inducedness is not required, and loops and pairs of parallel
edges are allowed as cycles in this auxiliary parameter.

Put $R=F-e$, where $F$ and $e$ are as in Section~2. The endpoints
$a=r_0^0$ and $b=\ell_1^0$ are the two attachment vertices of $R$.
For a selected set in $R$, every selected nonattachment vertex must
have internal selected degree two. At an attachment use state $0$
for unselected, state $1$ for selected with internal degree two, and
state $2$ for selected with internal degree one. State $2$ requires
a selected neighbour across the future attachment edge.

\begin{lemma}[Exact attachment table]\label{lem:edge-table}
The only feasible pairs of attachment states in $R$ are
$(0,0),(0,1),(1,0),(2,2)$. For each of these four pairs, the maximum
number of selected vertices is six.
\end{lemma}
\begin{proof}
An odd number of state-$2$ attachments is impossible by the degree
sum formula. For states with no $2$, the selected set is good in
$F-e$, giving the upper bound six by Lemma~\ref{lem:cap}.
For $(2,2)$, adding $e$ gives a good set in $F$, with the same bound.

State $(1,1)$ is impossible. Selecting $a$ forces both $\ell_0$
vertices to be selected; selecting $b$ similarly forces both $r_1$
vertices. Neither $Q_0$ nor $Q_1$ can be full, since the square
analysis in Lemma~\ref{lem:cap} would then empty the other squares.
Thus they are exactly these three-sets. Their doubled sides require
both sides of $Q_2$ to be selected. This makes $Q_2$ full, again
contradicting that analysis.

For explicit attainments, number $\ell_i^0,\ell_i^1,r_i^0,r_i^1$
as $4i,4i+1,4i+2,4i+3$, respectively, so $(a,b)=(2,4)$.
The following table gives selected six-sets:
\[
\begin{array}{c|l}
 (0,0)&1,3,5,7,9,11\\
 (0,1)&4,6,7,8,9,10\\
 (1,0)&0,1,2,8,10,11\\
 (2,2)&0,2,4,6,8,10.
\end{array}
\]
Their internal degrees verify the stated attachment states.
\end{proof}

For each edge $uv$ of $B$, take a fresh copy of $R$, remove $uv$,
and join $u$ to its attachment $a$ and $v$ to its attachment $b$.
For a loop, both joins use the same base vertex but distinct
attachment vertices. Call the result $\mathcal I(B)$.

\begin{theorem}[Cycle-packing encoding]\label{thm:inflation}
If $B$ is a cubic multigraph on $n$ vertices, then $\mathcal I(B)$
is simple and cubic, and
\[
 |V(\mathcal I(B))|=19n,\qquad
 \cind(\mathcal I(B))=9n+\nu_2(B).
\]
\end{theorem}
\begin{proof}
There are $3n/2$ edge copies, each adding twelve vertices. Every
attachment receives exactly one new edge, and every base vertex
retains its three incidences. Distinct copies have disjoint vertices,
so even loops and parallel edges in $B$ create no loops or parallel
edges in $\mathcal I(B)$.

In a good set of $\mathcal I(B)$, mark a base edge precisely when
its copy of $R$ has state $(2,2)$. Lemma~\ref{lem:edge-table} shows
that these are exactly the copies whose attachment edges occur in
selected cycles. Each selected base vertex has exactly two marked
incidences; an unselected base vertex has none. The marked edges
therefore form an ordinary $2$-regular subgraph of $B$, covering at
most $\nu_2(B)$ base vertices. Each edge copy contributes at most
six, proving the upper bound $6(3n/2)+\nu_2(B)$.

Conversely, choose a maximum ordinary cycle packing in $B$ and
select its base vertices. Use the displayed $(2,2)$ six-set in each
marked edge copy and the $(0,0)$ six-set in every other copy.
All selected vertices have degree two, and the resulting good set
has the asserted cardinality.
\end{proof}

For completeness, the elementary bound needed to apply this formula
also holds for multigraphs with loops.

\begin{lemma}\label{lem:ordinary-packing}
Every connected cubic multigraph $B$ on $n$ vertices satisfies
$\nu_2(B)\geq n/2+1$. If $B$ is bridgeless, then $\nu_2(B)=n$.
\end{lemma}
\begin{proof}
Choose a maximum ordinary cycle packing, covering a set $S$ of
$s$ vertices. Its complement $T$ induces a forest, since any cycle
there could be added to the packing. Let $c$ be the number of
components of this forest, and let $q$ count the edges of $B[S]$
outside the packing. Thus $|E(B[S])|=s+q$, the cut has $s-2q$
edges, and $|E(B[T])|=n-s-c$. Counting all $3n/2$ edges gives
\[
 3n/2=(s+q)+(s-2q)+(n-s-c),\qquad s=n/2+q+c.
\]
If $T$ is nonempty, then $c\geq1$. Otherwise $s=n$ and
$q=n/2\geq1$. This proves the first bound.

For a bridgeless cubic multigraph, a loop would force its remaining
incident edge to be a bridge. The graph is therefore loopless, and
Petersen's perfect-matching theorem~\cite{Petersen} gives a spanning
$2$-factor, proving $\nu_2(B)=n$.
\end{proof}

Consequently every connected $B$ satisfies
\[
 \cind(\mathcal I(B))\geq\frac{|V(\mathcal I(B))|}{2}+1.
\]
For bridgeless $B$, the exact value is
$\cind(\mathcal I(B))=10|V(\mathcal I(B))|/19$.
The additive constant one is sharp: take a tree with $t$ vertices
of degree three and $t+2$ leaves, and put a loop at each leaf.
Only those loops can occur in an ordinary cycle packing, so
$n=2t+2$ and $\nu_2(B)=t+2$. The resulting simple cubic graph has
order $38t+38$ and induced cycles vertex number $19t+20$.

\section{A criterion for nonuniform edge replacements}
The preceding exact table is unusually simple. The following
sufficient condition applies to different fragments on different
edges and explains why a larger collection of replacements cannot
produce a counterexample over a bridgeless base.

Let $B$ be a loopless bridgeless cubic multigraph on $n$ vertices.
For each edge $e$, let $R_e$ be a simple graph with two designated
degree-two attachments and all other degrees three. Replace $e$
by $R_e$ as before, and call the resulting simple cubic graph $G$.
For each feasible attachment state $(i,j)$, define
\[
 f_e(i,j)=\max_S\{2|S|-|V(R_e)|\},
\]
where the maximum is over locally valid selected sets in that state;
put $f_e(i,j)=-\infty$ for an infeasible state. Use the states of
Lemma~\ref{lem:edge-table}, and write
\[
 C_e=\max_{i,j\in\{0,1\}}f_e(i,j),\qquad
 D_e=f_e(0,0),\qquad A_e=f_e(2,2).
\]
Both $C_e$ and $D_e$ are finite, since the empty set is allowed.

\begin{theorem}[Matching criterion]\label{thm:matching-criterion}
Suppose every $A_e$ is finite. If
\begin{equation}\label{eq:local-criterion}
 \sum_{e\in E(B)}(3C_e+D_e+2A_e)\geq0,
\end{equation}
then $\cind(G)\geq |V(G)|/2$. In particular, it suffices that
$3C_e+D_e+2A_e\geq0$ for each edge separately.
\end{theorem}
\begin{proof}
First leave every base vertex unselected. In each fragment use
a closed state attaining $C_e$. These choices are compatible and
give a good set with doubled surplus
\[
 Z_0=\sum_e C_e-n.
\]
Next, for a perfect matching $M$ of $B$, select every base vertex.
Use state $(0,0)$ in fragments on $M$ and state $(2,2)$ in all
other fragments. Each base vertex then has exactly two selected
neighbours, and all other selected degrees are two by the local
states. This good set has doubled surplus
\[
 Z_M=n+\sum_{e\in M}D_e+\sum_{e\notin M}A_e.
\]

There is a probability distribution on perfect matchings of $B$
under which each edge occurs with probability $1/3$. Indeed, the
constant edge vector $1/3$ has degree sum one at every vertex.
Every odd vertex set has an odd cut, and bridgelessness makes that
cut have at least three edges. The vector therefore satisfies the
perfect-matching polytope inequalities; Edmonds' theorem
\cite{Edmonds} expresses it as a convex combination of perfect
matchings. Consequently,
\[
 \mathbb E Z_M=n+\frac13\sum_e(D_e+2A_e),\qquad
 \frac{Z_0+\mathbb E Z_M}{2}
 =\frac16\sum_e(3C_e+D_e+2A_e)\geq0.
\]
At least one of the constructed good sets has nonnegative doubled
surplus, which proves the claim.
\end{proof}

If $B$ has a proper three-edge-colouring, its three colour classes
provide the required distribution directly, so this case uses no
polyhedral result. Of the 27 retained two-attachment profiles
described below, 21 have a finite $(2,2)$ state. Independent subset
enumeration verifies $3C+D+2A\geq0$ for each of these 21 fragments.
The theorem therefore applies to every choice and orientation of
these fragments on every loopless bridgeless cubic base, with no
bound on its order. This is an application to a specified fragment
library; we have not proved the local inequality for arbitrary
fragments, and the hypothesis on the base excludes bridges.

\section{Two defects and a finite bipartite census}
The following reduction isolates a restricted class of possible
counterexamples. A set $D$ is called \emph{$1$-limited dominating}
if it dominates its complement and every vertex of $D$ has at most
one neighbour outside $D$. In a cubic graph the latter condition
means that $G[D]$ has minimum degree at least two. This parameter and its relation
to induced cycles are studied in~\cite{RRT}.

\begin{lemma}[Two-defect reduction]\label{lem:two-defect}
Let $G$ be cubic of order $n$, and let $D$ be a $1$-limited
dominating set of size $n/2+1$. Either $D$ is good, or its complement
$S$ is good and $G[D]$ has exactly two vertices of degree three.
In the second case, $G[D]$ consists of a theta graph or two cycles
joined by a path, together with zero or more cycle components.
Here a theta graph consists of three internally disjoint paths
with common endpoints.
\end{lemma}
\begin{proof}
Write $t=|D|$, $s=|S|=t-2$, and let $q$ count the vertices of
degree three in $G[D]$. Its other vertices have degree two, so
the cut has $t-q$ edges. Domination gives $t-q\geq s$, whence
$q\leq2$. The handshaking lemma makes $q$ even.
If $q=0$, $D$ is good. If $q=2$, the cut has exactly $s$ edges,
so every vertex of $S$ has exactly one neighbour in $D$ and
$G[S]$ is $2$-regular.

The two degree-three vertices belong to the same component,
since every component has an even number of odd-degree vertices.
Suppress all degree-two vertices in that component, temporarily
allowing loops and parallel edges. The resulting cubic multigraph
on two vertices has either three parallel edges, or one edge
joining the vertices and a loop at each. Expanding the suppressed
paths gives the asserted theta or joined-cycle structure.
Every remaining component is a cycle.
\end{proof}

For a bipartite $G$ in the second case, the two degree-three
vertices lie in opposite bipartition classes. To see this, write
$t_L,t_R$ for the numbers of vertices of $D$ in the two classes
and $a,b$ for their numbers of degree-three vertices in $G[D]$.
Counting internal edges in the two classes gives
$2t_L+a=2t_R+b$. The good set $S$ is a union of even cycles, so
its two classes have equal size. The cut is a matching saturating
$S$ and all the degree-two vertices of $G[D]$, giving
$t_L-a=t_R-b$. These equations imply $a=b=1$ and $t_L=t_R$.
Consequently all three theta paths have odd length; in the other
case the two cycles are even and the joining path has odd length.
All additional cycle components are even.

\begin{proposition}[A parity restriction]\label{prop:parity}
If $n\equiv0\pmod4$, a cubic bipartite graph of order $n$ has no
$1$-limited dominating set of size $n/2+1$. If $n\equiv2\pmod4$,
it has no such set of size $n/2$.
\end{proposition}
\begin{proof}
In the first case, both $n/2+1$ and $n/2-1$ are odd, whereas
Lemma~\ref{lem:two-defect} makes one of the corresponding sets a
union of even cycles. In the second case, a dominating set $D$
of size $n/2$ has at most $|D|$ cut edges and must dominate
$n/2$ outside vertices. Equality forces every vertex of $D$
to have internal degree two. Its odd cardinality is impossible
in a bipartite graph.
\end{proof}

We exhaustively generated the second case of
Lemma~\ref{lem:two-defect} for bipartite graphs of orders
$10,14,18,22,26$. A template specifies the even cycle lengths
of $S$, the core lengths just described, and the even cycle
lengths of the remaining components of $D$. For each template
we enumerate all colour-respecting matchings between $S$ and
the degree-two vertices of $D$. For a single $S$-cycle, rotation
allows one specified attachment to be fixed; no such reduction
is made when $S$ has several components. If $|S|=2h$, the numbers
of tested matchings per template are respectively
$(h-1)!h!$ and $(h!)^2$.

Fixing a bipartition for each template does not omit a gluing:
a global interchange fixes the core's colour convention, and
rotation by one vertex interchanges the colours of any even
cycle component. Repeated isomorphic completions are permitted.
The counts therefore refer to labelled completions, not distinct
isomorphism classes.

\begin{center}
\begin{tabular}{rrr}
\toprule
Order & Templates & Tested completions\\
\midrule
10 & 1 & 2\\
14 & 3 & 36\\
18 & 12 & 4,320\\
22 & 24 & 207,360\\
26 & 88 & 36,115,200\\
\bottomrule
\end{tabular}
\end{center}

Every completion admitted a good set of size at least $n/2+1$.
The search constructs an explicit set and directly checks its
cardinality and every selected degree before accepting a case.
Witnesses are checked during the run; the saved output consists
of aggregate counts and rolling digests, not a separate certificate
for every completion. Thus this is a reproducible finite computation.

The program \texttt{verify\_two\_defect\_search.py} independently
checks the structural list against all 22 complements of order
14 generated by \texttt{geng -qb -d2 -D3 14 15:15}, and all four
cycle factors of order 12 generated with degrees fixed at two.
It also compares the new counting solver against direct subset
enumeration on all 112 connected cubic graphs through order 12.
All checks pass. The source programs, templates, outputs and
scope statements are retained in the accompanying research record.

As a consequence of this computation and the preceding lemmas,
every cubic bipartite graph of order at most 26 that has a
$1$-limited dominating set of size at most $n/2+1$ satisfies
\eqref{eq:conj}. Indeed, counting the cut first gives $|D|\geq n/2$;
at equality $D$ itself is good. A set of size $n/2+1$ is handled
by the parity restriction, the good-set alternative, or the
enumeration. At order six the second alternative cannot occur,
since it would require a good set of order two.
There is no proof here that an arbitrary cubic graph has such a
dominating set, nor does the finite census establish the two-defect
case at unbounded order.

\section{Obstructions to exchange arguments}
An exchange argument would begin with a good set $S$ and replace
it by a larger good set. Define its \emph{deletion radius}
\[
 \rho_G(S)=\min\{|S\setminus T|:T\text{ is good and }|T|>|S|\},
\]
with value $+\infty$ if no improvement exists. The next statements
identify limitations of two possible exchange strategies. They
are not counterexamples to the induced cycles conjecture.

\begin{lemma}[Girth forces large exchanges]\label{lem:exchange-girth}
Let $G$ be a simple cubic graph of girth $g$. Let $S$ be good,
suppose its edge boundary is a matching, and let $q$ count the
vertices outside $S$ with no neighbour in $S$. Every good set
$T\ne S$ satisfies
\[
 g\leq6|S\setminus T|+q.
\]
In particular, $\rho_G(S)\geq\lceil(g-q)/6\rceil$.
\end{lemma}
\begin{proof}
Put $R=S\setminus T$, $A=T\setminus S$, and let $B$ be the
vertices of $S\cap T$ whose two selected incident edges change.
Write $r=|R|$, $a=|A|$, and $b=|B|$. Each vertex of $B$ loses
one neighbour in $R$ and gains one in $A$: two distinct pairs
of edges at a cubic vertex have exactly one common edge.
There are at most $2r$ edges from $R$ to $S\setminus R$, so
$b\leq2r$.

Every vertex of $A$, except possibly the $q$ vertices counted
in the hypothesis, has a neighbour in $S$. That neighbour must
belong to $R\cup B$. A retained vertex outside $B$ already has
its two selected neighbours and cannot acquire another.
Since the boundary is a matching, these neighbours in $R\cup B$
are distinct. Hence $a\leq r+b+q$.

The symmetric difference of the edge sets of $G[S]$ and $G[T]$
has degree two at precisely the vertices of $R\cup A\cup B$,
and degree zero elsewhere. It is nonempty because $S\ne T$,
so it contains a cycle. Therefore
\[
 g\leq r+a+b\leq2r+2b+q\leq6r+q.
\]
\end{proof}

\begin{theorem}[No uniform bounded exchange]\label{thm:unbounded-exchange}
For every nonnegative integer $r$, there is a connected simple
cubic bipartite graph $G$ with a good set $S$ of size
$|V(G)|/2-1$ such that $G-S$ has minimum degree two and every
distinct good set $T$ satisfies $|S\setminus T|>r$.
\end{theorem}
\begin{proof}
We first construct such two-defect graphs with arbitrarily
large girth. Fix an integer $g\geq3$ and put $k=2^g$.
Take a cycle $S$ of length $6k$, and a disjoint theta graph $D$
whose three paths each have length $2k+1$. Thus $|D|=6k+2$;
its two branch vertices have degree three and all others have
degree two. Its branch vertices are in opposite bipartition
classes. Each class of $S$ and each class of the degree-two
vertices of $D$ has size $3k$. Join these vertices by arbitrary
colour-respecting perfect matchings. Every resulting graph is
simple, cubic, bipartite and connected, with $S$ good and exactly
two defects outside $S$.

Among these finitely many completions, choose one minimizing
lexicographically the numbers of cycles of lengths
$3,4,\ldots,g-1$. Suppose its girth is $\ell<g$. The fixed
cycle and theta have girths $6k$ and $4k+2$, respectively, so a
shortest cycle uses a matching edge $uv$, with $u\in S$.
The union of the radius-$(g-1)$ balls around $u$ and $v$ has
at most
\[
 2\bigl(1+3(2^{g-1}-1)\bigr)=3\cdot2^g-4
\]
vertices. There are $3k=3\cdot2^g$ matching edges in the same
colour class as $uv$. Since these edges form a matching, one
of them, say $xy$ with $x\in S$, has both endpoints outside
those balls. Replace $uv,xy$ by $uy,xv$.

This switch preserves the template, simplicity and bipartiteness.
A new cycle using just one new edge has length greater than
$g$, since its remaining path joins the two separated endpoint
sets. A new cycle using both new edges splits into two old
paths. If these paths join $u$ to $x$ and $v$ to $y$, both
have length at least $g$. Otherwise they join $u$ to $v$ and
$x$ to $y$, avoiding the deleted edges; each has length at
least $\ell-1$, by the old girth. In this case the new cycle
has length at least $2\ell>\ell$.
Thus no cycle of length at most $\ell$ is created, while a
cycle of length $\ell$ is destroyed. This contradicts the
lexicographic choice. The chosen completion has girth at least $g$.

Now take $g=6r+3$. Lemma~\ref{lem:exchange-girth}, with $q=2$,
gives $6r+3\leq6|S\setminus T|+2$ for every distinct good set
$T$, and hence $|S\setminus T|>r$.
\end{proof}

The theorem disproves the assertion that a fixed number of deletions
always suffices to improve a two-defect set. It does not decide whether
a larger good set exists in every constructed graph. The distinction
matters: large girth alone supplies an obstruction to the proposed
local step, not to the original inequality.

There is also a barrier to a different move with no length bound.
Call a replacement of a good set $S$ by a good set $T$ a
\emph{single-cycle exchange} when
$E(G[S])\mathbin{\triangle}E(G[T])$ is the edge set of one
simple cycle. Every difference of two good sets in a cubic graph
is a disjoint union of cycles, but performing those cycles one
at a time need not preserve inducedness or the objective value.

\begin{theorem}[Traps for single-cycle improvement]\label{thm:cycle-trap}
For every integer $t\geq2$, there is a connected simple cubic
graph $G_t$ of order $16t+26$ and a good set $S_t$ such that
\[
 |S_t|=7t+14,\qquad \cind(G_t)=8t+13,
\]
and no sequence of nondecreasing single-cycle exchanges starting
at $S_t$ reaches a larger good set. For $t=2$, the exact loss
barrier is two: every single-cycle-exchange sequence reaching a
set larger than 28 visits a set of size at most 26, and a sequence
with sizes $28,27,26,29$ exists.
\end{theorem}
\begin{proof}
Use the sharp bridge construction of Section~5, taking the
internal vertices of its underlying tree to form a path on
$t$ vertices. There are $t$ triangles and $t+2$ copies of the
cap $P$. Leave every triangle unselected and choose the
seven-vertex root-containing set in every cap. The resulting
good set $S_t$ has size $7(t+2)$, whereas the exact optimum
from Section~5 is $8t+13$. Their difference is $t-1>0$.

Every edge joining different pieces is a bridge. Thus a simple
cycle, and hence a single-cycle exchange, lies wholly within
one cap or one triangle. Each cap already contributes its
maximum of seven vertices. Changing its root from selected to
unselected would leave at most six, so a nondecreasing sequence
cannot make that change. Every triangle has at least one
adjacent leaf cap, because the internal tree is a path.
Selecting the triangle while that cap's root is selected would
put both ends of a bridge in a good set, which is impossible.
Neutral cap moves therefore cannot unlock any triangle, and
no nondecreasing sequence reaches an improvement.

For $t=2$, each triangle is adjacent to two caps. An improvement
must select a triangle, since the four caps together contribute
at most 28. Immediately before the first triangle is selected,
both adjacent cap roots must be unselected. The total size at
that point is at most $6+6+7+7=26$.
This bound is attainable. With the numbering of Section~2,
replace the root-containing set
$\{0,2,4,6,8,10,12\}$ in either adjacent cap by
$\{0,3,5,6,8,10\}$. The latter is good and excludes the root;
the edge-set difference is a single cycle of length seven.
Doing this in both adjacent caps and then selecting their
triangle gives the stated sizes $28,27,26,29$.
\end{proof}

The accompanying verifier checks all 107 simple cycles of $P$
and thus all 430 simple cycles of the 58-vertex example. Exactly
88 flips yield good sets; none improves the initial size 28.
The escape sequence is checked directly. It also independently
verifies minimum deletion radii for 600 seeded two-defect
instances of orders 26, 38 and 50, exhausting 48,480 candidates
with smaller deletion counts. Those sampled radii are one or two;
their small values do not extend to the unbounded-girth family.

\section{Odd spanning factors and exact edge profiles}
An \emph{odd spanning factor} of a cubic graph $G$ is a spanning
subgraph $M$ in which each vertex has degree one or three. An
\emph{isolated edge} means a component isomorphic to $K_2$.
The following elementary reformulation
separates parity from the inducedness condition.

\begin{proposition}\label{prop:odd-factor}
Good sets of a cubic graph $G$ are in bijection with odd spanning
factors having no isolated-edge component. The correspondence is
\[
 S\longmapsto M_S=(V(G),E(G)\setminus E(G[S])),
 \qquad M\longmapsto L(M),
\]
where $L(M)$ is the set of degree-one vertices of $M$.
If $n=|V(G)|$, $k(M)$ is the number of components, and
$\beta(M)=|E(M)|-n+k(M)$, then
\begin{equation}\label{eq:factor-surplus}
 |L(M)|-\frac n2=n-|E(M)|=k(M)-\beta(M).
\end{equation}
Consequently, the conjecture is equivalent to the existence of an
odd spanning factor with no isolated edge and at most $n$ edges.
\end{proposition}
\begin{proof}
For a good $S$, vertices in $S$ have degree one in $M_S$, and all
others have degree three. An edge joining two degree-one vertices
in $M_S$ would be an edge of $G[S]$, contrary to its deletion.
Conversely, every edge outside an odd factor $M$ has both endpoints
in $L(M)$, since a degree-three vertex uses all its incident edges.
Each leaf therefore has two neighbours in $L(M)$ outside $M$.
An additional edge inside $L(M)$ would be an isolated edge of $M$.
In its absence, $G[L(M)]$ is $2$-regular, and the two maps are inverse.
Writing $\ell=|L(M)|$, the degree sum gives
$2|E(M)|=\ell+3(n-\ell)=3n-2\ell$; this proves the identities.
\end{proof}

\begin{corollary}\label{cor:odd-pseudoforest}
If a cubic graph has an odd spanning pseudoforest with no isolated
edge, it satisfies the half bound. More precisely, the leaves form
a good set of size $n/2+t$, where $t$ is the number of tree components.
\end{corollary}
\begin{proof}
A tree component contributes $1$ to $k-\beta$, and a unicyclic
component contributes zero. Apply \eqref{eq:factor-surplus}.
\end{proof}

This sufficient condition is precisely the earlier pseudoforest
complement condition in another form: $M_S$ consists of $G-S$ with
one pendant edge for each vertex of $S$. It is not a proved existence
theorem for arbitrary cubic graphs. Scott's perfect-forest theorem,
for which Gutin gives a short proof \cite{Gutin}, supplies an odd
spanning forest in every connected graph of even order, but permits
isolated-edge components. It therefore does not supply the missing
condition. Moreover, replacing ``pseudoforest'' by ``forest'' cannot
work universally: the graph $F$ of Section~2 has $\cind(F)=6$, whereas
an isolated-edge-free odd forest on its twelve vertices would give
at least seven leaves by Corollary~\ref{cor:odd-pseudoforest}.

\begin{proposition}\label{prop:factor-enumeration}
Let $G$ be connected and cubic of order $n$, and let $\mathcal O(G)$
be its odd spanning factors. Then $|\mathcal O(G)|=2^{n/2+1}$.
For $M\in\mathcal O(G)$ write $J(M)$ for its isolated-edge components,
regarded as a set of edges. For each edge $e=uv$, put
\[
 B_e=\cind(G-e),\qquad
 T_e=\max\{|S|:S\text{ good in }G,\ u,v\in S\},
\]
with $T_e=-\infty$ when no such set exists. Then
\begin{align*}
 B_e&=\max\{|L(M)|:M\in\mathcal O(G),\ e\in E(M),\ J(M)\subseteq\{e\}\},\\
 T_e&=\max\{|L(M)|:M\in\mathcal O(G),\ e\notin E(M),\ J(M)=\varnothing\}.
\end{align*}
All these optima and the complete good-set size distribution can be
computed in $O(n2^{n/2+1})$ time and polynomial working space.
\end{proposition}
\begin{proof}
Symmetric difference with $E(G)$ identifies odd factors with the
binary cycle space. A spanning tree gives a fundamental cycle basis
of dimension $|E(G)|-n+1=n/2+1$. Thus all factors are obtained exactly
once by toggling subsets of this basis.

For any odd factor, a leaf has exactly two neighbours joined to it
by edges outside the factor; both are leaves. The only extra edges
of $G[L(M)]$ are therefore the edges in $J(M)$. Accordingly,
$(G-e)[L(M)]$ is $2$-regular exactly when $e\in E(M)$ and
$J(M)\subseteq\{e\}$. Conversely, every good set in $G-e$ gives such
an odd factor by removing its induced edges from $G$. This proves
the formula for $B_e$. If $J(M)$ is empty, the two endpoints of an
edge $e$ are selected together exactly when $e\notin E(M)$; this
proves the second formula.

For each factor, compute its leaves and isolated edges in $O(n)$
time. If there are no isolated edges, update the histogram and all
eligible edge optima. If there is exactly one, update the deletion
optimum for that edge alone. Factors with two or more contribute
to none of these quantities. A binary Gray order toggles one
fundamental cycle per step; its length is at most $n$. The basis,
current factor, histogram, optima, and witnesses require only
polynomial space.
\end{proof}

Subdividing $e$ gives a rooted fragment of order $n+1$ with exact
root states $A_e=1+T_e$ and $B_e$. By Lemma~\ref{lem:double}, joining
two copies at their roots gives a counterexample precisely when
\[
 2B_e<n+1\quad\text{and}\quad B_e+T_e<n.
\]
An impossible root-in state is omitted from the maximum.
Thus the enumeration checks a complete family of rooted
counterexample tests for each input graph.

The implementation \texttt{odd\_factor\_enumerate.cpp} and its
standard-library verifier use different search spaces. Direct
enumeration of all vertex subsets on all 112 connected cubic graphs
through order twelve agrees in every histogram coefficient and all
1,899 edge profiles; it checks 369,040 subsets. A further seeded
sample contains 240 labelled graphs of orders 24 through 42, including
sharp constructions, one to eight degree-preserving switches from
them, retained search candidates, and random cubic graphs. Exhausting
261,726,208 odd factors gives 11,526 exact edge profiles. Every input
satisfies the original half bound, and every deletion has $B_e\geq n/2$;
none passes the counterexample test above. Five inputs attain the
original half bound. All maximum and edge-profile witnesses are
retained and checked directly. These 240 inputs are a sample, not an
exhaustive census up to order 42. Their upper bounds rely on complete
cycle-space enumeration; they are not asserted on the strength of
witnesses alone.

\section{Repairing odd pseudoforests}
For an odd spanning factor $M$, write $j(M)=|J(M)|$ for its number
of isolated-edge components. A \emph{cycle exchange} replaces its
edge set by $E(M)\mathbin{\triangle}E(C)$, where $C$ is a simple
cycle of $G$. The resulting factor is still odd. Here we permit
intermediate factors with isolated edges, but require every
intermediate factor to be a pseudoforest. This is a different
state space from the good-set exchanges in Section~9.

The following first step has an unrestricted proof.

\begin{lemma}\label{lem:nonalternating}
If $P$ is a perfect matching of a nonempty finite simple cubic graph $G$,
then $G$ has an induced cycle that is not $P$-alternating.
\end{lemma}
\begin{proof}
If $G$ is not bipartite, a shortest odd cycle is induced and is
not alternating. Suppose instead that $G$ has bipartition $(A,B)$.
Orient matching edges from $A$ to $B$, and all other edges from
$B$ to $A$. A cycle is directed exactly when it is $P$-alternating.
If every induced cycle were alternating, this would be a cyclic
orientation of $G$ in the sense of Speyer~\cite{Speyer}.
Speyer's decomposition theorem implies that cyclically orientable
graphs are $2$-degenerate: their blocks are built from edges and
cycles by successively attaching a cycle along one edge.
In particular, each block is series-parallel, and the block
decomposition preserves $2$-degeneracy. Thus a nonempty such graph
has a vertex of degree at most two, contrary to cubicity.
\end{proof}

\begin{theorem}[A repair from every perfect matching]\label{thm:matching-repair}
Every perfect matching $P$ of a nonempty finite simple cubic graph admits
a cycle exchange to an odd spanning pseudoforest $M$ with
$j(M)<j(P)=n/2$.
\end{theorem}
\begin{proof}
Choose an induced nonalternating cycle $C$ by
Lemma~\ref{lem:nonalternating}. Put $r=|V(C)|$,
$k=|E(C)\cap P|$, and $b=r-2k>0$. The $b$ vertices of $C$
not covered by $E(C)\cap P$ have distinct matching partners
outside $C$, because $C$ is induced. Toggle $E(C)$.

If $k=0$, the changed component is $C$ with one pendant edge
at every vertex. It is unicyclic. If $k>0$, deleting the $k$
matching edges from $C$ leaves disjoint paths; adding the
unchanged matching edges at their $b$ internal vertices attaches
distinct pendant leaves. These components are trees. Unaffected
components remain isolated edges. Thus the new factor is an odd
spanning pseudoforest.

Exactly $k+b$ old matching components are affected. When $k>0$,
let $a$ be the number of remaining cycle paths consisting of a
single edge; these are precisely the new isolated edges, and
$a\leq k$. Set $a=0$ when $k=0$. Consequently
\[
 j(M)=\frac n2-k-b+a\leq\frac n2-b<\frac n2.\qedhere
\]
\end{proof}

This theorem establishes only the first step. After the exchange,
nontrivial tree or unicyclic components may be present, so the
matching hypothesis is lost. Reapplying the theorem is not justified.
A proposed sufficient repair assertion for the full conjecture was:
\begin{quote}
From every odd spanning pseudoforest with $j>0$, there is a finite
sequence of cycle exchanges through odd spanning pseudoforests
such that $j$ never increases and eventually strictly decreases.
\end{quote}
Every connected cubic graph has even order and an odd spanning
forest. Iterating the stated assertion would reach $j=0$ and then
prove the half bound by Corollary~\ref{cor:odd-pseudoforest}.
The assertion is \emph{false}, as proved in
Theorem~\ref{thm:pseudo-barrier} below. We retain the earlier finite
tests and first-step obstructions to make clear why they did not
settle this stronger assertion.

\begin{proposition}[Certified exchange obstructions]\label{prop:repair-traps}
The following finite examples exist.
\begin{enumerate}
\item A connected cubic graph of order 14 has an odd spanning
forest with one isolated edge from which no single cycle exchange
both preserves the pseudoforest condition and reduces $j$.
One neutral exchange followed by an improving exchange does succeed.
\item A connected cubic graph of order 16 has an odd spanning
pseudoforest with 15 edges and one isolated edge, but every
nontrivial cycle exchange preserving the pseudoforest condition
has at most 14 edges.
\item The one-step obstruction in the first assertion also occurs
in a connected bipartite cubic graph of order 36. Again, two
exchanges suffice, the first neutral in $j$.
\end{enumerate}
None of these graphs is a counterexample to the induced cycles
conjecture.
\end{proposition}
\begin{proof}[Finite certificate and verification]
Number edges in graph6 order: increasing larger endpoint, then
increasing smaller endpoint, starting at index zero. An integer
mask selects exactly the edges with set bits. The first two
instances have the following graph6 codes and factor masks:
\[
\begin{array}{ccl}
14&\texttt{M?`@?aGqBCCoDOCW?}&1296836,\\
16&\texttt{O???C@\_cSoAgQ\_P\_D\_@D?}&4716667.
\end{array}
\]
The first factor is three copies of $K_{1,3}$ and one $K_2$.
The second has components with $(|V|,|E|)$ equal to
$(6,6),(8,8),(2,1)$.

For the third graph, take the bipartite double cover of the
18-vertex graph with graph6 code
\texttt{Q????A?oE?kCEGL?DO?i?Ao?F??}.
Label its two sheets $0,\ldots,17$ and $18,\ldots,35$; replace
each edge $uv$ by $u(v+18)$ and $(u+18)v$. Use factor mask
$2812704630600655$ in the resulting graph6 edge order.
Its components have orders $10,6,6,12,2$; only the order-12
component is unicyclic, and the others are trees.

Direct vertex-path enumeration gives respectively 93, 196 and
11,976 simple cycles. Each cycle is toggled, and connected
components of the resulting factor are inspected directly.
In the 14-vertex case, the complete table of allowed moves is
\[
\begin{array}{c|rrrrrr}
(|E(M')|,j(M'))&(7,7)&(8,5)&(9,3)&(10,1)&(11,2)&(12,1)\\
\text{number of cycles}&4&20&32&15&1&2.
\end{array}
\]
In particular $j$ never drops below one in a single move.
A two-move repair is the sequence of factor masks
$1296836,313647,1930178$.
In the 16-vertex case, the numbers of allowed moves grouped by
new edge count $8,9,10,11,12,13,14$ are
$2,14,32,37,35,33,13$; there are none with at least 15 edges.
In the 36-vertex case, every allowed move has $j\geq1$;
exactly 153 have $j=1$. A two-move repair has masks
\[
2812704630600655,\quad
2812730224997562,\quad
2812704640340943.
\]
The last factors in the two repair sequences have no isolated
edge and have respectively nine and 24 leaves, giving good sets.
Direct subset enumeration gives $\cind=10$ for the 16-vertex graph.
Thus all three satisfy the original half bound.

The standard-library program
\texttt{verify\_pseudoforest\_repair.py} reproduces the cycle
lists, full move tables, and shortest repairs. Its output retains
all graph edges and factors explicitly, so the certificates do
not require graph6 software. Component traversal and direct
edge-parity recursion are independent of the C++ disjoint-set
and cycle-space implementation. These are computer-checked
finite statements, not an unrestricted repair theorem.
\end{proof}

The stronger assertion allowing neutral moves passed the following
finite tests. Every odd spanning factor was enumerated for each
input; for each pseudoforest with $j>0$, all direct moves were
examined until an improvement was found, or a breadth-first search
explored moves with unchanged $j$ until an improvement was reached.
\begin{center}
\small
\begin{tabular}{lrrr}
\toprule
Input collection & Graphs & Odd factors & \shortstack{Factors needing\\a neutral move}\\
\midrule
All connected cubic graphs, $n\leq18$ &45,982&44,513,544&214\\
All connected cubic graphs, $n=20$ &510,489&1,045,481,472&1,234\\
Seeded sample, $24\leq n\leq42$ &240&261,726,208&144\\
Bipartite double covers, $28\leq n\leq36$ &39&18,776,064&48\\
\bottomrule
\end{tabular}
\end{center}
All tested factors needing neutral moves admitted a subsequent
strict improvement. In the complete census through order twenty,
each such factor needs exactly two exchanges: one neutral and one
improving. At order twenty there are 920,048,096 odd pseudoforests,
of which 889,866,061 contain isolated edges. The 1,234 factors
requiring a neutral move occur in 291 distinct graphs.
Completeness of the graph inputs up to isomorphism relies on
\texttt{geng -c -d3 -D3 20}; the factor search exhausts each graph's
cycle space. These counts concern the stronger repair assertion;
the original conjecture was already checked through order twenty-two
in the literature~\cite{KSS}.
The last two collections are samples, not
censuses of their orders. The double covers come from the 39
distinct graphs supporting the 214 small neutral-move cases.
Their lifted initial factors all admit immediate improvements;
the 48 neutral-move cases are different factors found by exhausting
each cover's cycle space. The 36-vertex certificate above is one
of these. A further 7,616 completions of its fixed factor, obtained
by degree-preserving switches among the other edges, produced
ten one-step obstructions, all of vertex connectivity two.
Three-connectivity does not remove the obstruction, as the next
finite certificate shows.

\begin{proposition}[A three-connected one-step obstruction]\label{prop:three-connected-trap}
There is a three-connected cubic graph $G$ of order twenty and an
odd spanning pseudoforest $M$ with one isolated edge such that no
single cycle exchange preserving the pseudoforest condition
reduces $j(M)$. Two exchanges suffice. Nevertheless,
$\cind(G)=14> |V(G)|/2$.
\end{proposition}
\begin{proof}[Finite certificate and verification]
In the graph6 and edge-mask convention above, take
\[
 G=\texttt{S??????o@\_@\_gCGoWCAD?Cg?X??WO?IG?},\qquad
 M=983379956.
\]
Deleting any one or two vertices leaves the graph connected.
The factor $M$ consists of the isolated edge $\{1,14\}$ and an
18-vertex unicyclic component, whose unique cycle is the triangle
on $\{5,11,19\}$. Direct vertex-path enumeration gives 664 simple
cycles. Toggling each and inspecting the resulting components
gives no pseudoforest with $j=0$, and exactly 57 with $j=1$.
For the neutral moves, the new edge counts $15,16,17,18,19$ occur
respectively $9,18,19,9,2$ times. The complete move table is retained
with the certificate.

A shortest repair has factor masks
\[
 983379956\ \longrightarrow\ 983451959\
 \longrightarrow\ 977257976.
\]
Their respective pairs $(|E(M)|,j(M))$ are $(19,1),(17,1),(19,0)$.
The last factor's leaves form the good set
\[
 \{0,1,2,4,7,9,10,13,14,17,18\}.
\]
Independent enumeration of all $2^{20}$ vertex subsets gives the
good-set size polynomial
\begin{align*}
 \sum_{S\text{ good}}x^{|S|}
 ={}&1+2x^3+4x^4+2x^5+5x^6+9x^7+10x^8\\
 &+13x^9+13x^{10}+14x^{11}+11x^{12}+6x^{13}+x^{14}.
\end{align*}
Thus the maximum is fourteen. The standard-library verifier
\texttt{verify\_round9.py} checks the connectivity, all cycles,
the repair, and this polynomial independently of the C++ census.
\end{proof}

The census finds exactly three three-connected graphs supporting
one-step obstructions at order twenty, with one such factor in each.
The other 1,231 factors have vertex connectivity one or two
(152 and 1,079 factors, respectively). Of the 1,234 factors, 1,218
have $j=1$ and sixteen have $j=2$; all supporting graphs have
girth three. None of the 214 neutral
cases through order eighteen is three-connected. Hence twenty is
the minimum order of a three-connected one-step obstruction within
the complete census. This is a computational minimality statement,
not a general theorem about higher girth or larger orders.
The independent verifier checks every one of the 1,448 reported
neutral cases through order twenty, including the absence of a
direct improvement and an explicit two-exchange repair. It does
not independently repeat the entire 556,471-graph census.

A different shortcut is to minimize the global potential
$h(M)=|E(M)|+j(M)$ over all odd spanning factors. For an odd
pseudoforest, $h(M)\leq n$: indeed $|E(M)|=n-t(M)$, and every
isolated edge is a tree component, so $j(M)\leq t(M)$.
However, a minimum of $h$ need not be attainable with $j=0$.
For the ten-vertex graph with graph6 code \texttt{I?BeeOwM?},
exhausting its 64 odd factors gives
\[
 \min_{j(M)>0}h(M)=8<9=\min_{j(M)=0}h(M).
\]
A factor with mask 9700 attains the left side and consists of
two copies of $K_{1,3}$ and one $K_2$. Direct enumeration of all
$2^{10}$ vertex subsets gives $\cind=6$, so this example also
satisfies the original conjecture. It rules out this particular
global minimization argument and the stronger rounding claim
$\cind(G)\geq |L(M)|-j(M)$, whose right side here is seven.

Independently, direct edge-parity enumeration checks all 12,296
odd factors in the 112 connected cubic graphs through order 12.
Of these, 10,911 are pseudoforests, and all 9,334 having an
isolated edge admit an immediate improvement. It also checks
17,356 induced-cycle repairs from all 1,098 perfect matchings in
these graphs against Theorem~\ref{thm:matching-repair}.
The finite results do not bridge the gap from the proved first
step to the unrestricted repair assertion, which is disproved in
Theorem~\ref{thm:pseudo-barrier}.

\section{Single-edge replacement and repair simulation}
The fragment $R=F-e$ from Section~6 provides a useful test of
attempts to enlarge the repair obstructions. Its effect can be
determined for every cubic base graph, without a size bound.
Let $H$ be a simple cubic graph, let $uv\in E(H)$, and form
$H_R$ by deleting $uv$, adding a disjoint copy of $R$, and joining
$u$ and $v$ to its ports $a=2$ and $b=4$, respectively.

\begin{theorem}[Exact single-edge replacement]\label{thm:single-replacement}
With $B_{uv}=\cind(H-uv)$ and $T_{uv}$ as in Section~10,
\[
 \cind(H_R)=6+\max\{B_{uv},T_{uv}\}
             =6+\max\{\cind(H-uv),\cind(H)\}.
\]
In particular, replacement by $R$ preserves the conjectured half
bound whenever it holds for $H$.
\end{theorem}
\begin{proof}
Consider any good set in $H_R$. By Lemma~\ref{lem:edge-table},
its restriction to $R$ has at most six vertices and one of the
boundary states $(0,0),(0,1),(1,0),(2,2)$.
In any of the first three states, each selected vertex outside
$R$ has no selected neighbor in $R$: a port in state $1$ is
selected but already has two selected neighbors inside $R$, so
its outside neighbor must be unselected. The outside restriction
is therefore good in $H-uv$ and has size at most $B_{uv}$.
In state $(2,2)$ both $u$ and $v$ must be selected. Replacing
their edges to $R$ by $uv$ makes the outside restriction good in
$H$, with both ends of $uv$ selected. Its size is at most $T_{uv}$.
This proves the upper bound.

Conversely, a largest good set of $H-uv$ extends by a six-vertex
$(0,0)$ state in $R$. A good set counted by $T_{uv}$ extends by
a six-vertex $(2,2)$ state. Thus the upper bound is attained.
Every good set in $H$ either contains both ends of $uv$, or
remains good after $uv$ is deleted. Since $T_{uv}\leq\cind(H)$,
this also proves the second equality. Finally,
$|V(H_R)|=|V(H)|+12$, so the increase of at least six preserves
the half bound.\qedhere
\end{proof}

Consequently a smallest counterexample cannot contain an induced
copy of $R$ whose two outside neighbors are distinct and nonadjacent:
smoothing that copy to one edge would give a smaller simple cubic
graph to which Theorem~\ref{thm:single-replacement} applies.
The restrictions on the outside neighbors ensure that smoothing
creates neither a loop nor a parallel edge.

We next show that a natural choice of factor inside $R$ also
preserves repair paths. All vertex labels below are those of $F$
in Section~2. Define a spanning tree $T$ of $R$ by
\begin{align*}
 E(T)=\{&(0,2),(1,2),(1,3),(1,11),(4,6),(4,7),\\
        &(5,6),(6,8),(8,10),(8,11),(9,11)\}.
\end{align*}
The ports $2,4$ have degree two in $T$; every other vertex has
degree one or three. Let
\[
 P=(2,0,3,5,7,4),\qquad E(U)=E(T)\mathbin{\triangle}E(P).
\]
Then $U$ is the six-cycle $(1,3,5,6,8,11,1)$ with one leaf
attached at each cycle vertex. Its leaves are
$0,2,4,7,9,10$. In particular, $U$ is an odd spanning unicyclic
graph with no isolated edge.

For an odd spanning factor $M$ of $H$, define its lift by
\[
 \widehat M=\begin{cases}
 (M-uv)\cup T\cup\{ua,vb\},&uv\in E(M),\\
 M\cup U,&uv\notin E(M).
 \end{cases}
\]
Here copies of the fragment are understood to be disjoint from
$V(H)$.

\begin{theorem}[A lift preserving repair paths]\label{thm:repair-lift}
If $M$ is an odd spanning pseudoforest of $H$, then $\widehat M$
is an odd spanning pseudoforest of $H_R$, with
\[
 j(\widehat M)=j(M)-\mathbf1_{\{uv\in J(M)\}},\qquad
 |L(\widehat M)|=|L(M)|+6.
\]
Every simple-cycle exchange in $H$ lifts to one simple-cycle
exchange in $H_R$. Suppose $uv\notin J(M_0)$ and $M_0$ has a
repair of length $\ell$ through odd pseudoforests with $j$
nonincreasing. Then $\widehat M_0$ has such a repair of length
at most $\ell$.
\end{theorem}
\begin{proof}
If $uv\in E(M)$, its replacement by the tree $T$ and the two
attachment edges preserves the cycle rank and connectedness of
the component containing $uv$. All old vertex degrees are
unchanged; the ports have degree three, and other new vertices
have degree one or three. No new isolated edge is created.
An old isolated edge disappears precisely when it is $uv$.
If $uv\notin E(M)$, the new factor adds the separate unicyclic
component $U$, and all old components are unchanged. Both
constructions add exactly six leaves. This proves the displayed
identities and the pseudoforest assertion.

Let $C$ be a simple cycle in $H$. If $uv\notin E(C)$, use the
same cycle in $H_R$. Otherwise concatenate the edge $ua$, the
internal path $P$, and the edge $bv$ in place of $uv$. The
resulting cycle has edge set
\[
 E(\widehat C)=(E(C)\setminus\{uv\})\cup E(P)\cup\{ua,bv\}.
\]
It is a simple cycle. The identity $T\triangle U=P$ gives
\[
 \widehat{M\triangle C}=\widehat M\triangle\widehat C.
\]
For a sequence through odd pseudoforests, all intermediate lifted
factors are therefore pseudoforests.

Truncate the original repair at its first strict decrease, so
$j(M_i)=j(M_0)$ before the last state. Initially $uv\notin J(M_0)$,
and hence $j(\widehat M_0)=j(M_0)$. Along the lifted sequence,
either $uv$ first becomes isolated in some $M_i$, immediately
giving $j(\widehat M_i)=j(M_0)-1$, or the strict decrease in
$j(M_i)$ is reached first. Before that first decrease the lifted
values remain constant. Taking this prefix gives a repair of
length at most $\ell$.
\end{proof}

This theorem rules out a particular search strategy: inserting
this factor pair along an initially nonisolated edge cannot turn
a two-exchange repair into one requiring three exchanges. It
does not say that every odd pseudoforest of $H_R$ is such a lift,
and supplies no repair when none is known in $H$.

The standard-library verifier \texttt{verify\_round10.py} checks
the explicit fragment, the factor identities, and the cycle lifts.
It also checks the exact optimum formula by vertex-subset
enumeration on all 24 single-edge replacements of connected
cubic graphs through order six. A separate exploratory test
retains 2,895 triangle expansions of the three order-20
obstructions, at orders 22 through 42, and 216 replacements by
$R$ at order 32. All admit repair witnesses of one or two moves;
every retained witness is independently checked. These samples
are not exhaustive at their orders.

The 216 factors occur in 54 distinct labelled graphs. Exhausting
all 7,077,888 odd factors of these inputs gives 5,895,198
pseudoforests, of which 5,878,454 have isolated edges. Every one
of the latter admits a repair within two exchanges; 1,838 require
a neutral first move. Explicit repair paths for all 1,838 are
independently verified. For one such factor in each of the 54
graphs, the verifier also independently exhausts every simple
cycle, totaling 968,802 cycles, and confirms that two exchanges
are necessary. This does not independently repeat the full factor
census. The proof above, rather than these finite tests,
establishes the unrestricted lifting statement.

\section{Unbounded barriers to pseudoforest repair}
The monotone repair assertion fails in graphs with bridges, despite
all the finite checks in Section~11. This failure can be arbitrarily
large. In this section, a \emph{repair sequence} for $M$ is a finite
sequence $M=M_0,M_1,\ldots,M_s$ of odd spanning pseudoforests in
which consecutive factors differ by a simple cycle and
$j(M_s)<j(M)$. Intermediate defect counts are allowed to rise.
Define the \emph{excess repair barrier} by
\[
 h_G(M)=\min_{(M_i)}\bigl(\max_i j(M_i)-j(M)\bigr),
\]
with value $+\infty$ if there is no repair sequence. A monotone
repair must have $h_G(M)=0$; positive barriers rule it out.

Use the thirteen-vertex cap $P$ from Section~2, with
$\ell_i^j=4i+j$, $r_i^j=4i+2+j$, and root $r=12$.
Two of its good sets are
\[
 X=\{0,3,5,6,8,10\},\qquad
 Y=\{0,2,4,6,8,10,12\}.
\]
Write $N_X=E(P)\setminus E(P[X])$ and define $N_Y$ similarly.
The subgraph $N_X$ is connected and unicyclic, with root degree
two and all other degrees odd. The subgraph $N_Y$ has the root
isolated and a unicyclic component on the other twelve vertices.
Their edge-set difference is the seven-cycle
\begin{equation}\label{eq:cap-switch}
 Q=(0,2,12,4,6,5,3,0).
\end{equation}
These descriptions follow directly from the edge list in Section~2;
neither subgraph has an isolated-edge component.

For $k\geq3$, construct $B_k$ from two disjoint $k$-cycles by
joining one designated vertex of each with a bridge. Attach a
separate copy of $P$, by its root, to every other vertex of the
two cycles. Thus there are $2k-2$ caps. The resulting graph is
connected, simple and cubic, with
\[
 n_k=|V(B_k)|=2k+13(2k-2)=28k-26.
\]
Define $M_k$ to contain every bridge and a copy of $N_X$ in every
cap, and no edge of either $k$-cycle. It has one isolated edge,
namely the bridge between the designated vertices. Its other
$2k-2$ components are unicyclic graphs of order fourteen, each
obtained from $N_X$ by adding its adjacent cycle vertex as a leaf.
In particular, $M_k$ is an odd spanning pseudoforest with $j(M_k)=1$.

\begin{lemma}[Cost of freeing a cap]\label{lem:free-cap}
Let $N$ be a spanning pseudoforest of $P$ in which the root has
degree zero or two and every other vertex has degree one or three.
Let $j_P(N)$ count its isolated-edge components and put
\[
 w(N)=j_P(N)+\mathbf1_{\{\deg_N(r)=0\}}.
\]
If the component of $N$ containing the root is a tree, then
$w(N)\geq1$.
\end{lemma}
\begin{proof}
For root degree zero the second summand already equals one.
Suppose the root has degree two and $j_P(N)=0$. Let $L$ be the
set of nonroot vertices of degree one in $N$. Every edge outside
$N$ has both ends in $L$, because the root and all degree-three
vertices use every edge available to them in $P$. Each vertex
of $L$ consequently has two neighbours in $L$ outside $N$.
There is no additional edge of $N$ inside $L$, since it would be
an isolated-edge component. Hence $P[L]$ is $2$-regular and
avoids the root, so $|L|\leq6$ by Lemma~\ref{lem:cap}.

The degree sum gives $|E(N)|=19-|L|$. If the root component
were a tree, the pseudoforest would have at most twelve edges,
giving $|L|\geq7$, a contradiction.
\end{proof}

\begin{theorem}[Exact unbounded repair barriers]\label{thm:pseudo-barrier}
For every $k\geq3$,
\[
 \cind(B_k)=\frac{n_k}{2},\qquad h_{B_k}(M_k)=k-1.
\]
The shortest repair sequence for $M_k$, allowing increases in
$j$, has exactly $k$ exchanges. One such sequence has defect counts
\[
 1,2,\ldots,k,0.
\]
Thus there is no monotone repair sequence, and no constant bounds
the necessary temporary increase in $j$, even for starting factors
with just one isolated edge.
\end{theorem}
\begin{proof}
Every simple cycle of $B_k$ lies in one cap or is one of the
two displayed $k$-cycles, since every edge between pieces is a
bridge. In any good set, each displayed cycle is therefore either
entirely selected or entirely unselected. Both cannot be selected,
because their joining bridge would then lie in the induced
$2$-regular subgraph. If neither is selected, the caps contribute
at most $14(k-1)$ vertices. If one is selected, its adjacent cap
roots must be unselected. The cap bounds $A=7$, $B=6$ give
\[
 \cind(B_k)=\max\{14(k-1),\ k+6(k-1)+7(k-1)\}
            =14k-13=\frac{n_k}{2}.
\]
Both bounds are attainable by the specified cap states.

Consider any repair sequence starting at $M_k$. Until the first
exchange on a displayed $k$-cycle, all its vertices have factor
degree one and all bridges remain present. If $N_i$ are the
current restrictions to the caps, the isolated-edge count is
exactly
\begin{equation}\label{eq:barrier-cost}
 j=1+\sum_{i=1}^{2k-2} w(N_i).
\end{equation}
The initial isolated bridge persists throughout this period, so
some displayed cycle must be exchanged before $j$ can reach zero.

Immediately before that first exchange, each of the $k-1$ cap
components attached to that cycle must be acyclic. Otherwise
the exchange adds the displayed cycle to a component that already
contains a cap cycle, violating the pseudoforest condition.
Lemma~\ref{lem:free-cap} and \eqref{eq:barrier-cost} now give
$j\geq1+(k-1)=k$. Moreover, each of those caps started with its
root in a cyclic component and so must have undergone at least
one exchange. Since one exchange cannot cross a bridge, at least
$k-1$ cap exchanges and the displayed-cycle exchange are necessary.

For attainment, choose one displayed cycle. In each of its
$k-1$ caps, toggle $Q$ from \eqref{eq:cap-switch}, changing
$N_X$ to $N_Y$. Each move leaves a unicyclic component of order
twelve and makes the cap's attachment edge an isolated edge.
Thus all intermediate factors are odd pseudoforests, with
$j=1,2,\ldots,k$. Now toggle the chosen displayed cycle.
It forms a unicyclic component with its $k$ adjacent factor
leaves and absorbs all $k$ isolated edges. Every remaining
component is unicyclic, so $j=0$. This uses exactly $k$ exchanges
and attains the stated barrier.
\end{proof}

The smallest member of this family has order 58. Its repair has
$j$-values $1,2,3,0$ and factor-edge counts $57,56,55,58$.
Its final good set has 29 vertices. It is a counterexample to the
monotone repair assertion, \emph{not} to the induced cycles
conjecture. The existence of an isolated-edge-free odd spanning
pseudoforest in every cubic graph remains unproved; refuting this
particular way of reaching one does not refute that existence claim.

The independent program \texttt{verify\_round11.py} exhausts all
$2^{19}$ cap-edge subsets and all $2^{13}$ cap-vertex subsets.
There are 128 cap parity states. Exactly twelve have $w=0$ and
are pseudoforests; their roots lie in cyclic components, and they
form a connected graph under cap-cycle exchanges. The program
checks all $12\cdot107$ moves from these states. For $B_3$, it
also checks all 430 simple cycles at $M_3$ and directly verifies
the $12^4=20,736$ global states in the neutral component, including
the 41,472 attempted exits obtained by flipping a displayed triangle.
Every such exit violates the pseudoforest condition. Cap moves
remain in this product or increase $j$, so this accounts for the
entire neutral component. Explicit attaining sequences and good
sets are checked for $3\leq k\leq12$, at orders 58 through 310.

These bridge obstructions should not be read as evidence of the
same barrier in bridgeless graphs. For the canonical bipartite
double covers of these ten examples, an independently checked
pair of simple-cycle exchanges gives $j$-values $2,1,0$ and
a good set of order $n/2+4$. Their orders range from 116 to 620;
connectedness and the absence of bridges are also checked directly.
These certificates concern the specified covers only.

\section{Checks and the remaining obstacle}
\subsection{Direct tests of the original bound}
For an edge $uv$ of a cubic graph $B$, an \emph{edge-triangle
replacement} deletes $uv$ and introduces three vertices $a,b,w$
and the five edges $ua,ab,bv,aw,bw$. It adds two branch vertices
and one degree-two vertex. We applied this operation to every
unordered pair of distinct original edges in each connected simple
cubic graph $B$ of orders twelve, fourteen and sixteen generated
by \texttt{geng}. The two edges may share an endpoint.
For the resulting graph $H$, writing $n=|V(B)|$, we have
\[
 |V(H)|=n+6,\qquad |E(H)|=\tfrac32n+8,
 \qquad |E(H)|-|V(H)|=\tfrac n2+2.
\]
Every construction admitted a good set containing at least
$n/2+2$ branch vertices. The counts refer to indexed constructions;
different entries may be isomorphic.
\begin{center}
\begin{tabular}{rrrr}
\toprule
Base order & Base graphs & Edge-pair constructions & Completion order\\
\midrule
12 & 85 & 13,005 & 44\\
14 & 509 & 106,890 & 46\\
16 & 4,060 & 1,120,560 & 48\\
\bottomrule
\end{tabular}
\end{center}
There are exactly two missing incidences in each $H$, so its cubic
completion has order $n+32$. The exact completion formula proves
the original half bound for all these completions.

The search enumerates the binary cycle space of $H$. In a subcubic
graph an Eulerian edge set has degree zero or two at every vertex,
so its support is a union of vertex-disjoint cycles. Checking whether
that support induces a $2$-regular graph discards precisely the
supports with additional chords or edges between cycles. Conversely,
every good set occurs as such a support. For a connected graph this
gives an exhaustive search with $2^{m-n+1}$ edge states, although the
implementation stops when a qualifying positive witness is found.
The separate standard-library program \texttt{verify\_round12.py}
reconstructs all 1,240,455 graphs and checks every saved vertex-set
certificate directly, without using the search algorithm or an
optimization package. The claimed completeness of the base list up
to isomorphism relies on \texttt{geng}.

A second search tested 1,595 distinct labelled cubic graphs obtained
by degree-preserving edge changes from the sharp family in Section~13
and from rings of the fragment $R$. Their orders range from 36 to 114.
Every graph has a directly verified half-bound witness, including one
case resolved by retrying a time-limited query. Fourteen named cubic
graphs, with orders between 8 and 126 and LCF descriptions taken
from the SageMath source \cite{Sage}, also have verified witnesses
strictly exceeding half their orders. These are lower bounds on
individual graphs, not certified maxima or exhaustive order censuses.
The verifier reconstructs each LCF graph independently.

Motivated by Theorem~\ref{thm:bridge-symmetry}, a further search made
20,000 degree-preserving edge-switch proposals on rooted fragments
of orders $13,25,29,37$, corresponding to symmetric cubic joins of
orders $26,50,58,74$. The proposals contain 14,088 distinct labelled
rooted graphs; they are not an exhaustive census. Every proposal
has a verified half-bound witness in its full doubled cubic graph.
Of the 20,000 records, 19,845 terminate once a strictly positive
surplus is certified; their recorded local values are lower bounds.
The remaining 155 records have exact profiles, all attaining the
half bound upon doubling. These comprise 95 distinct labelled
rooted graphs. The standard-library verifier
\texttt{verify\_round13.py} independently checks every full cubic
witness and rechecks all 95 exact profiles by vertex-assignment
recursion, using degree constraints to propagate forced choices.
The search itself uses cycle-space enumeration. No timeout or
early positive exit is interpreted as an optimum or a counterexample.

Lemma~\ref{lem:triangle-pair} is checked separately on all 3,615
edge replacements in labelled subcubic graphs of orders two through
five. All 7,984 good-set lifts have the claimed gain. The proof of
the lemma, rather than this finite check, establishes its unrestricted
statement. None of these searches or reductions proves
\eqref{eq:conj} for arbitrary cubic graphs.

\subsection{A strictness census and a further synthesis attempt}
A separate pass over all 556,471 connected cubic graphs of orders
four through twenty finds precisely one isomorphism class attaining
the half bound: the twelve-vertex Franklin graph $F$. For each of
the other 556,470 graphs, a saved good set has at least $n/2+1$
vertices. These are positive certificates, not claimed exact optima.
For the single equality case, exhaustive vertex-subset enumeration
confirms optimum six, twenty maximum good sets, and the fact that
every edge belongs to at least one maximum good set. An explicit
relabeling identifies the graph with $F$.

The verifier \texttt{verify\_round14.py} reconstructs every graph,
checks connectedness and cubic degree, verifies the saved set, and
compares the certificate list with the two existing \texttt{geng}
input lists row by row. As before, isomorphism coverage relies on
\texttt{geng}; the witness checks are independent of the search.
Proposition~\ref{prop:one-triangle} then proves the branch inequality
for every single edge-triangle replacement in these bases. A strict
base already supplies the needed extra vertex, and in $F$ the
through-edge optimum is six for every edge. There are 16,539,825
indexed choices of a base graph and an edge; this is a consequence
of the census and the proposition, rather than separate enumeration
of all those expanded graphs. Their cubic completions also satisfy
\eqref{eq:conj}, by Theorem~\ref{thm:completion}.

Finally, two bounded synthesis runs sought a twenty-five-vertex
rooted graph with $(A,B)\leq(12,12)$ or $(A,B)\leq(13,11)$,
respectively. Either profile would violate the half bound after
doubling. The runs fix the root neighbours and permit deletion of
at most six edges of a specified labelled seed, together with the
same number of insertions. They examined thirteen and ten candidates
before their wall-time limits. Independent cycle-space enumeration
checks both optima for all twenty-three candidates and verifies a
half-bound witness in every full fifty-vertex cubic double. Neither
run establishes infeasibility in its restricted search region.
The degree-propagation oracle and cycle-space checker are also
cross-checked against all vertex subsets in all 843 labelled
subcubic graphs through order five. The triangle formula is checked
by full vertex-subset enumeration for all 387 edge expansions of
cubic bases through order ten and of $F$.

\subsection{Certificates for simultaneous replacements}
Theorem~\ref{thm:cut-certificate} has exact rational certificates
for a list of 2,249 connected cubic bases through order twenty and
one twenty-four-vertex base. The former list was selected by a
cycle-space search for bases with optimum at most $n/2+2$; the
certificate verification does not depend on that selection criterion.
For 1,477 of the 2,250 bases, a uniform distribution on the listed
good sets suffices. For the other 773, linear programming
proposed a distribution, which was converted to rational coefficients.
The separate verifier checks every listed set, the sum and sign of
the coefficients, and \eqref{eq:cut-certificate} using exact
\texttt{Fraction} arithmetic. Numerical tolerances are not acceptance
criteria.

Each certificate covers every subset of edges simultaneously,
and Corollary~\ref{cor:triangle-chains} extends it to arbitrarily
long chains. Thus these are certificates for infinite families of
cubic graphs, rather than a claim about a fixed list of expanded
orders. The twenty-six-vertex sharp bridge graph does not meet
the criterion, consistently with the upper bound proved after
Theorem~\ref{thm:cut-certificate}; its recorded distribution attains
$25/2$ exactly.

The implementation of \eqref{eq:simultaneous} is independently
cross-checked against all edge subsets of the eight connected
cubic bases through order eight. Its optimum formula is then
checked against every vertex subset of 551 expanded graphs:
all marking subsets for $K_4$, and all marking subsets of size
at most two for the remaining bases through order eight. The
verifier constructs and degree-checks a half-bound witness in
all 551 full cubic completions, whose orders are at most one hundred.

An earlier search in this continuation followed sequences of
replacements that decrease or preserve branch surplus, to depth
four. It produced 11,193 distinct labelled expanded graphs; every
saved lower-bound witness is checked directly. This restricted
search does not cover sequences that temporarily increase surplus.
A separate simultaneous-marking constraint model reported
infeasibility for 1,387 selected small bases before its wall-time
limit and for two further tight bases of orders twenty-four and
twenty-six. No formal solver proof files are retained; these
statuses are exploratory records. The exact rational distributions,
not these solver statuses, justify the infinite-family assertions.
All new data and checks are recorded by \texttt{verify\_round15.py}.
The original conjecture remains unresolved.

\subsection{Normal form and conditional certificates}\label{sec:round16}
The reductions in Theorem~\ref{thm:simple-base} settle the loop
and parallel-edge alternatives left by
Lemma~\ref{lem:triangle-pair}. They are structural proofs for
arbitrary orders. A standard-library verifier checks the
neighbourhood identity and its lifts on all 12,911 labelled
subcubic graphs on one through six vertices. It tests 17,925
nonempty good sets and 17,945 lifts. The short parallel-edge
configuration is checked in 2,463 bounded outside environments,
with 4,850 good-set lifts, covering equal, adjacent, and nonadjacent
outside neighbours. A further 55 loop and parallel-chain
configurations check the neighbourhood reduction with longer chains.
These finite checks support the implementation of the proofs; they
are not an enumeration of all prospective counterexamples.

For the sharp twenty-six-vertex double of $P$, two conditional
cut certificates have ten rational terms each and cut lower bound
thirteen. The verifier reconstructs every Eulerian edge set,
support, chord set, and cut; checks positive coefficients summing
to one; and checks \eqref{eq:conditional-cut} using exact fractions.
The two cases partition all $2^{39}=549{,}755{,}813{,}888$ indexed
zero-or-one marking choices. This is a certificate argument, not
individual enumeration. Corollary~\ref{cor:triangle-chains} then
covers arbitrary chain lengths on this base, which the original
unconditional cut criterion could not certify.

The four-set root certificates are checked directly on $P$.
Both root states are also checked arithmetically for all $2^{19}$
marking patterns. For the 191 patterns with at most two marked
edges, actual rooted graphs and full cubic doubles are constructed;
every doubled graph has a directly checked half-bound witness.
Their orders are at most ninety. These are positive witnesses,
not claimed optima for the expanded graphs.

The program \texttt{verify\_round16.py} performs these checks
without an optimization package; it saves counts and data hashes.
The data directory retains the rational certificates. The original
conjecture remains unresolved, including the unexpanded cubic case
in Theorem~\ref{thm:simple-base}.

\subsection{Marked-forest restrictions and rooted search}\label{sec:round17}
The proofs of Proposition~\ref{prop:marked-sparsity} and
Lemma~\ref{lem:marked-quotient} hold without an order bound.
For an implementation check, \texttt{verify\_round17.py}
compares the full subset inequalities with the contracted-forest
criterion for all 21,568 markings of the eight connected cubic
bases through order eight. Of these markings, 13,284 are forests
and 4,276 satisfy the stronger sparsity condition. For markings
of size at most three, the verifier constructs 10,351 neighbourhood
deletions and checks all 45,390 resulting good-set lifts. These
counts describe finite checks of the formulas, not a resolution
of the remaining counterexample search.

Two further constraint searches sought rooted graphs $Q$ of order
29 with $A(Q),B(Q)\leq14$. Such a graph would yield a 58-vertex
counterexample by doubling. The searches fix the two root neighbours
and bound the number of deleted seed edges by six or four, with
degree-preserving insertions. The radius-six run stopped on a query
timeout after one candidate; the radius-four run reached its wall
time limit after six candidates. Neither region is proved
infeasible. Independently enumerating binary cycle-space states
checks the exact rooted profiles of all seven candidates, and
the verifier checks a half-bound witness on each full cubic double.
None is a counterexample. Search records, exact verification counts,
and hashes are retained with the source.

\subsection{Direct descent and four further certificates}\label{sec:round18}
A direct search on cubic graphs ordered candidates first by their
maximum good-set size and then by the number of good sets attaining
that size. It used degree-preserving edge switches, occasional
upward moves, and restarts. Accepted moves were evaluated exactly
by binary cycle-space enumeration; early exits recorded only
positive lower-bound witnesses. Twelve seeds, including vertex
expansions and sharp graphs assembled from the established
attachments, led to 30,000 switch proposals at orders
$24,26,28,30,32,36,38,42$. Together with 185 retained vertex
expansions, the data contain 30,185 records and 28,733 distinct
labelled graphs. All retained witnesses are checked directly.
An independent vertex-assignment recursion verifies the optima
and maximum-set counts of all 45 recorded best graphs.
No counterexample was found. This is a bounded heuristic search,
with no claim of exhaustive coverage at these orders.

Four selected graphs were then tested for simultaneous edge-triangle
replacements. A constraint solver reported infeasibility for all
four counterexample queries; no formal infeasibility proofs were
retained. Independently of these statuses, explicit rational cut
certificates establish stronger positive results:
\begin{center}
\begin{tabular}{rrrr}
\toprule
Base order & Certificate terms & $\tau$ & $\lceil\tau-n/2\rceil$\\
\midrule
26 & 84 & $96/7$ & 1\\
28 & 10 & $29/2$ & 1\\
30 & 10 & $31/2$ & 1\\
32 & 12 & $33/2$ & 1\\
\bottomrule
\end{tabular}
\end{center}
The first certificate is uniform on maximum good sets; the others
use rational distributions proposed by linear programming.
The verifier checks every selected degree, coefficient, probability
sum, and cut total using exact fractions. The quantitative form of
Theorem~\ref{thm:cut-certificate} and the chain reduction show that
every cubic completion of arbitrary triangle-chain replacements on
each of these four bases has induced cycles vertex number at least
half its order plus one. These are proved infinite-family bounds,
not a proof of the unrestricted conjecture.

The programs \texttt{cubic\_multiplicity\_descent.cpp} and
\texttt{verify\_round18.py}, seed construction source, compressed
search records, rational certificates, and exact verification hashes
are retained with the manuscript. Non-best search records are
independently certified only as positive witnesses, rather than
independent optimum computations.

\subsection{Exhaustive edge neighbourhoods of two sharp graphs}
\label{sec:round20}
A direct counterexample search starts from the sharp 24- and 26-vertex
seeds retained in \texttt{round20\_data/seed24.g6} and
\texttt{seed26.g6}. The forbidden property is a good set of size at
least $n/2$; no auxiliary inequality is substituted for the original
conjecture. A constraint-generation search adds a clause excluding
every qualifying good set of each candidate. Its two runs produced
nine and seven candidates and stopped at wall-time limits, with
1,496 and 1,616 accumulated clauses. These solver outcomes are
inconclusive, not unsatisfiability certificates.

A separate enumeration completely checks the specified finite regions.
For each seed $G$, it generates every labelled simple cubic $H$ on
the same vertices satisfying $|E(G)\setminus E(H)|\leq4$.
For each set of $k$ deleted seed edges, pair its $2k$ endpoint
incidences in every possible way. Reject loops, repeated new edges,
and any new edge already belonging to the seed. This generates
exactly the desired edge changes: cubicity ensures that the added
edges pair precisely the endpoint incidences of the deleted edges.
Different pairings can yield the same labelled graph when incidences
share a vertex, so the following counts are reconnection records,
not counts of distinct graphs. Disconnected candidates are included.
\begin{center}
\begin{tabular}{rrr}
\toprule
Deleted edges $k$ & 24-vertex seed & 26-vertex seed\\
\midrule
0 & 1 & 1\\
1 & 0 & 0\\
2 & 984 & 1,182\\
3 & 35,592 & 47,464\\
4 & 1,937,772 & 2,849,236\\
\midrule
Total & 1,974,349 & 2,897,883\\
\bottomrule
\end{tabular}
\end{center}
Every record has a good set of size at least $n/2$. Most reuse a
previously found witness; exact vertex-assignment search supplies
383 and 967 additional witnesses, respectively. Thus neither of
these entire finite neighbourhoods contains a counterexample.
This conclusion does not exclude graphs outside the two regions.

The implementation is \texttt{exploratory/edge\_ball\_enumerate.cpp}.
The independent standard-library verifier \texttt{verify\_round20.py}
checks all 1,350 additional witnesses directly. It separately
enumerates the cycle spaces of all sixteen constraint-search
candidates and confirms their reported optima and qualifying-set
counts. For a test of the reconnection generator itself, an independent
degree-sequence recursion enumerates all 70 labelled cubic graphs
on six vertices and all 19,355 on eight vertices. For one seed of
each order, the radius-four regions agree exactly with the C++
generator: 54 and 1,147 distinct graphs, respectively. These tests
validate the small cases and retained witnesses; they do not constitute
an independent re-enumeration of the two large regions.

\subsection{Earlier checks and limitations}
The accompanying standard-library program
\texttt{verify\_bridge\_reduction.py} enumerates the good sets of $F$,
$F-e$, and $P$ and checks their claimed maxima.
It constructs the completing witnesses for all 843 labelled
subcubic graphs on one through five vertices and verifies their
degrees, cardinalities, and the surplus identity.
It also checks thirteen sharp tree constructions, including orders
from 26 through 186. These finite checks support the implementation;
the proofs above establish the unrestricted statements.

The program \texttt{verify\_edge\_inflation.py} independently
enumerates all $2^{12}$ subsets of $R$ to verify the complete
attachment table. It checks ordinary packing optima for twelve
small auxiliary multigraphs, the corresponding induced-cycle
witnesses after inflation, and eleven members of the new family
attaining the additive constant one.

Further searches found no violation among 9,000 local changes to
tight graphs of orders 24, 26, and 42; every optimization in those
runs finished with a certified optimum from the solver.
A floating-point weighted search also found no violation among
39,866 connected bridgeless cubic graphs of order 18. The latter
result is exploratory and is not included in the earlier exact
rational certificate collection through order 16.

The branch formulation was also tested directly on all 4,597,428
connected simple graphs of minimum degree two and maximum degree
three on three through seventeen vertices, as generated by
\texttt{geng}. All satisfied \eqref{eq:branch}; 618 attained equality.
The subset search stops early only after finding a strict lower-bound
witness; a reported equality requires exhaustion of all subsets that
could improve the score. Completeness up to isomorphism relies on
the generator. A separate enumeration checked all 1,225,155
single-edge deletions from the 45,982 connected cubic graphs through
order eighteen. These too satisfied the branch inequality.
An additional 9,000 degree-preserving local changes to graphs with
two degree-two vertices produced no violation.

For three-port fragments, 73,352 vertex deletions from cubic graphs
through order sixteen were reduced by componentwise domination
and port permutations to 259 representative profiles. An independent
standard-library verifier re-enumerates all subsets of each retained
fragment and checks its entire boundary table. All 202,020 tested
two-fragment joins and 4,900 further heuristic assemblies satisfied
the conjectured bound. These computations do not exhaust arbitrary
assemblies or larger fragments.

A further enumeration covered all 1,059,165 connected simple graphs
on nineteen vertices with exactly one degree-two vertex and all
other degrees three. In the notation of Lemma~\ref{lem:double},
each admitted witnesses with $B\geq9$ and $\max(A,B)\geq10$.
These bounds rule out the doubling construction for those fragments.
The search stops once both witnesses are found and does not compute
every exact optimum. Separately, 1,225,155 edge deletions from cubic
graphs through order eighteen reduce to 27 retained two-attachment
tables under componentwise domination and reversal of attachments.
The program \texttt{verify\_boundary\_profiles.py} independently
enumerates every subset of each retained representative and verifies
its complete table and the local inequalities used in
Theorem~\ref{thm:matching-criterion}. The other six retained tables
have an infeasible $(2,2)$ state and are not covered by that theorem.
The separate program \texttt{verify\_matching\_criterion.py} checks
451 explicitly constructed good sets in 82 uniform and nonuniform
assemblies over $K_4$ and the Petersen graph, including the averaging
identity in the proof.

\begin{samepage}
The separate program \texttt{verify\_further\_reductions.py}
checks two failures of stronger proof strategies.
The three-connected ten-vertex graph with graph6 code
\texttt{ICOedO[X?} cannot be covered by two good sets.
For the sixteen-vertex graph with graph6 code
\texttt{O??CA?oBCDT?U?D\_CgAA\_}, four good sets can cover every
vertex at least twice, but every such four-set multiset uses some
edge at least three times. Thus imposing an additional edge-use
bound of two fails. Their maximum good-set orders are six and ten,
respectively, so neither is a counterexample to \eqref{eq:conj}.

Corollary~\ref{cor:equivalent} leaves the central inequality
\eqref{eq:branch} unproved. The reduction removes neither the
unbounded size nor the global compatibility difficulty.
No general proof or counterexample is claimed.
\end{samepage}

\section{Disjoint Petersen four-poles}
\label{sec:petersen-poles}
The recent counterexamples to Petersen coloring in~\cite{GJMMMU}
provide a separate source of cubic graphs to test. The following
argument rules out a broader class of potential counterexamples to
\eqref{eq:conj}. No claim of priority is made for this elementary
packing argument.

Let $W$ be the Petersen graph with two adjacent vertices deleted.
Use vertex set $\{2,3,4,5,6,7,8,9\}$ and edge set
\[
 \{23,27,34,38,49,57,58,68,69,79\}.
\]
Its four degree-two vertices, or ports, are $2,4,5,6$.
For every pair of ports there is an induced cycle whose intersection
with the ports is exactly that pair:
\begin{center}
\begin{tabular}{cc}
\toprule
Selected ports & Vertex set of the induced cycle\\
\midrule
$2,4$ & $2,3,4,7,9$\\
$2,5$ & $2,3,5,7,8$\\
$2,6$ & $2,3,6,7,8,9$\\
$4,5$ & $3,4,5,7,8,9$\\
$4,6$ & $3,4,6,8,9$\\
$5,6$ & $5,6,7,8,9$\\
\bottomrule
\end{tabular}
\end{center}
The entries are vertex sets, not cyclic orderings. Each has at least
five vertices, as can be checked directly from the ten edges of $W$.

\begin{theorem}\label{thm:petersen-packing}
Suppose a simple cubic graph $G$ contains $k$ pairwise vertex-disjoint
induced copies of $W$. Then $c_{\rm ind}(G)\geq5k$. In particular,
if $|V(G)|=8k+r$ and $r\leq2k$, then $G$ satisfies
\eqref{eq:conj}.
\end{theorem}
\begin{proof}
For each copy create two demand vertices. On the other side of a
bipartite graph place one resource vertex for each edge of $G$
leaving at least one copy. Join both demands of a copy to its four
leaving edges. These edges are distinct because the copy is induced.
For a set of $s$ copies, its $4s$ port incidences use at least $2s$
distinct resources: an edge can meet at most two copies.
For any set of demand vertices, let $s$ be the number of represented
copies. Its neighbourhood therefore has at least $2s$ resources,
whereas it contains at most $2s$ demands. Hall's condition holds.

Choose a matching saturating all demands. In each copy select the
cycle in the table corresponding to its two assigned ports. Select
no other vertices of $G$. No leaving edge has both endpoints
selected, since that would assign its resource to two copies.
Thus each selected vertex has exactly its two neighbours on its
local cycle, and the union is induced and 2-regular. It has at least
$5k$ vertices. Finally, $5k\geq(8k+r)/2$ when $r\leq2k$.
\end{proof}

\paragraph{Application and independently checked certificates.}
The two 52-vertex examples linked in~\cite{GJMMMU}, House of Graphs
identifiers 57244 and 57278, each contain six disjoint induced copies
of $W$, leaving four other vertices. The theorem gives at least
30 selected vertices, exceeding the required 26.
The constructive matching implementation actually returns witnesses
of orders 31 and 30, respectively; neither number is asserted to be
an optimum. The graph data, six embeddings, and selected vertices
are retained in \texttt{round19\_data/petersen\_examples.json}.
The standard-library program \texttt{verify\_round19.py} independently
checks simplicity, cubicity, connectivity, every induced embedding,
and degree two at every selected vertex. It also enumerates all
256 subsets of $W$ to check the local cycle table.
The graph data were retrieved on October 28, 2026 from
\href{https://houseofgraphs.org/graphs/57244}{House of Graphs 57244}
and \href{https://houseofgraphs.org/graphs/57278}{57278}.

\paragraph{Why the obvious smoothing argument stops.}
Write $x_p$ for the outside neighbour of port $p$, temporarily
assuming these four neighbours are distinct. The natural smoothing
removes $W$ and adds $x_2x_6$ and $x_4x_5$, when this produces a simple
graph. To preserve a chosen outside set, a selected $x_p$ requires
port $p$ to have the selection status of its paired outside neighbour.
An exhaustive local check of all sixteen outside selection patterns
finds an extension with at least five internal vertices in the first
fifteen patterns, but no extension when all four outside vertices
are selected. The individual witnesses and exact local optima are
retained in \texttt{round19\_verification.json}.

The exceptional case has a short direct explanation. All four ports
must then be selected, with exactly one selected internal neighbour.
Writing $y_v\in\{0,1\}$ for selection of the four nonports gives
\[
 y_3+y_7=y_3+y_9=y_7+y_8=y_8+y_9=1.
\]
Hence $y_3=y_8$ and $y_7=y_9=1-y_3$. If $3,8$ are selected, vertex
3 has the three selected neighbours $2,4,8$. Otherwise $7,9$ are
selected, and vertex 7 has the three selected neighbours $2,5,9$.
Both cases are impossible. Thus this local smoothing is not an
unconditional extension lemma. The packing theorem does not require
such an extension, but also does not cover graphs with too few copies
of $W$. The unrestricted conjecture remains unresolved.

\section{Three-port certificates and vertex contraction}
\label{sec:three-port-contraction}
Let $R$ have odd order $q$, three designated degree-two ports, and
all other vertices of degree three. The fragment need not be connected.
A local selected set has degree two at each selected nonport. At each
port use state $0$ for unselected, $1$ for selected with internal
selected degree two, and $2$ for selected with internal selected
degree one. Let $f_R(s)$ be the maximum doubled surplus $2|S|-q$
for boundary state $s\in\{0,1,2\}^3$, with $-\infty$ for an
infeasible state. Across an external edge, compatible pairs are
exactly $(0,0),(0,1),(1,0),(2,2)$.

\begin{theorem}[A three-port contraction criterion]
\label{thm:three-port-contraction}
Suppose
\[
 f_R(0,0,0)\geq-1,
\]
and, for each choice of two ports, there is a state having $2$ at
those ports, either $0$ or $1$ at the third port, and surplus at least
$1$. Let a simple cubic graph $G$ contain $R$ as an induced fragment,
with its three external edges leading to distinct outside vertices.
Let $G_0$ replace $R$ by one vertex adjacent to those three vertices.
Then
\[
 c_{\rm ind}(G)\geq c_{\rm ind}(G_0)+(q-1)/2.
\]
Consequently a smallest counterexample to \eqref{eq:conj} cannot
contain such a fragment with $q>1$ and distinct outside neighbours.
\end{theorem}
\begin{proof}
The distinct-neighbour assumption makes $G_0$ simple and cubic.
Take a good set in $G_0$, and let $v$ be the replacement vertex.
If $v$ is unselected, use the $(0,0,0)$ witness in $R$, adding at
least $(q-1)/2$ vertices. Every outside selected vertex sees an
unselected port in place of the unselected $v$.

If $v$ is selected, exactly two of its neighbours are selected.
Use the corresponding witness with state $2$ at those two ports.
The third outside neighbour is unselected, so either state $0$ or
state $1$ at its port is compatible. The local witness has at least
$(q+1)/2$ vertices; replacing $v$ therefore gains at least
$(q-1)/2$. All selected degrees are two. Taking a maximum good set
of $G_0$ proves the inequality. The order also increases by $q-1$,
so the half-bound lifts from the smaller graph.
\end{proof}

The criterion needs only four feasible witnesses, rather than a
complete optimum table. It can also be applied successively to
substitutions at vertices of any cubic graph already satisfying the
conjecture.

\begin{theorem}[One-port choices in arbitrary assemblies]
\label{thm:one-port-assembly}
Consider disjoint fragments whose ports each receive one external
edge, with every external edge joining distinct fragments. Suppose
for each fragment either $f_R(0,0,0)\geq0$, or there is a port $p$
for which
\[
 f_R(0,0,0)+f_R(\mathbf e_p)\geq0,
\]
where $\mathbf e_p$ denotes state $1$ at $p$ and $0$ at the other
ports. Then the resulting cubic graph satisfies \eqref{eq:conj}.
\end{theorem}
\begin{proof}
In fragments with nonnegative all-zero surplus use that witness.
Each remaining fragment has two choices: the all-zero witness with
surplus $a_R$, and the one-port witness with surplus $b_R$, where
$a_R+b_R\geq0$. Two one-port choices conflict only when the two
chosen ports are the ends of the same external edge. Each fragment
has just one chosen port, so these conflicts form a matching.

For each conflicting pair choose one endpoint uniformly to use its
one-port witness and the other to use its all-zero witness. For
unpaired fragments choose either witness with probability one half.
Every realization is compatible: no external edge has two selected
ends. Each variable fragment has expected surplus
$(a_R+b_R)/2\geq0$, and the fixed fragments have nonnegative surplus.
Some realization has nonnegative total surplus. Its local good sets
therefore form an induced union of cycles on at least half the
vertices of the assembly.
\end{proof}

\paragraph{Larger fragments and certified joins.}
We computed all vertex-deletion profiles of 42 retained cubic graphs
of orders 24 through 32, giving 1,238 indexed fragments. The search
program \texttt{exploratory/three\_pole\_large.cpp} enumerates the
cycle space of the original graph. A local set with two degree-one
ports closes through the deleted vertex; a local set with no such
ports is already a union of cycles. These are the only possibilities
by the degree-sum formula. Testing induced degrees in the fragment
therefore recovers every feasible state, including paths that are
not good sets in the original cubic graph.

Of the 1,238 fragments, 1,190 satisfy
Theorem~\ref{thm:three-port-contraction}, and 1,230 satisfy
Theorem~\ref{thm:one-port-assembly}. The latter 1,230 can be mixed
arbitrarily, with arbitrary port orientations, in assemblies of any
size subject to that theorem's external-edge condition. The two
counts describe overlapping classes, not disjoint classifications.
Seven of the eight exceptions to the one-port criterion satisfy the
contraction criterion. None of these sufficient conditions is asserted
to be necessary.

For pairwise joins, canonically permuting ports leaves 594 distinct
tables. Componentwise domination reduces these to 45 retained lower
tables. Every ordered pair of retained tables was checked with all
six port matchings: 12,150 checks in total. In each case compatible
local witnesses have total doubled surplus at least zero; the
minimum certified surplus is zero. Domination and port relabelling
therefore certify the half-bound for every pairwise join from the
full 1,238-fragment library. These joined graphs have orders 46
through 62. This does not cover arbitrary assemblies of the eight
exceptional fragments.

The independent standard-library verifier \texttt{verify\_round21.py}
checks all 13,927 retained feasible-state witnesses directly, the two
sufficient criteria, the domination certificates, and all pairwise
compatibility checks. These are positive certificates: their validity
does not require the table entries to be optimal or the other states
to be infeasible. The verifier does not independently establish those
upper-bound or infeasibility claims. Data and witnesses are retained
in \texttt{round21\_data}, with verification results in
\texttt{round21\_verification.json}. The contraction criterion still
leaves graphs without these fragments, and the original conjecture
remains unresolved.

\section{The remaining disconnected fragment and mixed assemblies}
\label{sec:round22}
The unique indexed fragment failing both sufficient criteria of
Section~\ref{sec:three-port-contraction} is obtained by deleting
vertex 27 of its retained 28-vertex source. It is disconnected: its
components have orders 14 and 13, with two and one ports respectively.
The following exact calculation resolves arbitrary assemblies of
copies of this fragment, although mixed assemblies remain a separate
question.

Recall the two-port fragment $R=F-e$ from
Lemma~\ref{lem:edge-table}, with ports $u,v$. Form $R^+$ by adding
vertices $a,b$ and edges $ua,ab,bv$, making $a,b$ the new ports.
Then $R^+$ has fourteen vertices. Its only feasible boundary states
are
\[
\begin{array}{c|ccc}
 \text{state}&00&11&22\\\hline
 \text{maximum selected vertices}&6&8&8\\
 \text{doubled surplus}&-2&2&2.
\end{array}
\]
Indeed, in state $00$ use a maximum closed state in $R$, of order
six. In state $11$, both $u,v$ must be selected with internal degree
one in $R$, giving its state $22$ and six vertices, plus $a,b$.
In state $22$, the edge $ab$ supplies the internal neighbour of
each new port, so $u,v$ must be unselected; use state $00$ of $R$,
again with six vertices. A state $01$ or $10$ is impossible because
a saturated selected new port needs both its neighbours, including
the other new port. States with exactly one $2$ are excluded by
parity. These arguments and Lemma~\ref{lem:edge-table} establish
both upper bounds and attainments.

The other component is the thirteen-vertex cap $P$ of Section~2,
with root-out and root-in doubled surpluses $-1$ and $1$.
Its root cannot have local state $2$, again by parity. Explicit
isomorphisms from $R^+$ and $P$ to the two components of the
exceptional fragment are retained and independently checked.

\begin{theorem}[Exact cap-and-chain formula]\label{thm:cap-chain}
Join disjoint copies of $R^+$ and $P$ by pairing all their ports,
one new edge per port, to obtain a simple cubic graph $G$.
Let $r_{\rm cyc}$ be the number of $R^+$ copies belonging to
components of the fragment-incidence multigraph that contain no
$P$ copy. Then
\[
 c_{\rm ind}(G)=\frac{|V(G)|}{2}+r_{\rm cyc}.
\]
In particular every such assembly satisfies \eqref{eq:conj},
including every assembly made solely from copies of the exceptional
disconnected three-port fragment above.
\end{theorem}
\begin{proof}
In the incidence multigraph, $R^+$ copies have degree two and $P$
copies have degree one. Thus its components are cycles of $R^+$
copies or paths of $R^+$ copies terminated by two $P$ copies.
Parallel incidence edges are allowed; an incidence loop at $R^+$
would duplicate its existing edge $ab$ and is excluded by simplicity.

For a cyclic component with $k$ copies of $R^+$, use state $22$
in every copy. This gives doubled surplus $2k$, which is also an
upper bound since each local surplus is at most two.

Consider a path with $k\geq1$ copies of $R^+$. A $22$ state would
propagate along the path to a cap, which has no state $2$. Thus
each copy uses state $00$ or $11$. Put $x_i=0$ or $1$ accordingly.
Compatibility gives $x_i+x_{i+1}\leq1$. The maximum contribution
of each end cap is $1-2x_i$ for its adjacent end copy. The total
maximum doubled surplus for this state sequence is
\[
 \sum_{i=1}^k(4x_i-2)+2-2x_1-2x_k
 =2\sum_{i=1}^{k-1}(x_i+x_{i+1}-1)\leq0.
\]
An alternating sequence attains zero; for $k=1$ the sum is empty
and both choices attain zero. A path with no $R^+$ copy is just
two caps joined at their roots, whose maximum doubled surplus is
also zero. Local maximizing witnesses combine because their port
states are compatible. Adding the contributions over components
proves the formula.
\end{proof}

\paragraph{Checks and the remaining mixed case.}
The independent verifier \texttt{verify\_round22.py} enumerates all
24,576 vertex subsets of the two components to recover their exact
local tables. It checks the decomposition against the original
source graph, the two isomorphisms, and 200 independent transfer
calculations for paths through length 100 and cycles of lengths two
through 100. The theorem itself applies to unbounded lengths.

The separate bounded search \texttt{exploratory/exception\_assemblies.py}
uses all eight exceptional indexed fragments, allowing changes to
fragment choices and external port pairings. It retains 388 mixed
search records with four through twenty fragments. Another 56
records use uniform choices on fixed cubic bases. Full expanded
graph orders range from 100 to 580. For each assembly, a boundary
constraint model maximizes total doubled surplus. Of the 444 solver
results, 443 are reported optimal and one is feasible only. Every
record supplies a nonnegative-surplus witness, and each witness is
independently checked for degree two at every selected vertex in
the full simple cubic graph. No solver objective is interpreted as
an optimum on a merely feasible or unknown result.

All seven uniform assemblies of the disconnected exceptional
fragment attain the exact formula above, including the one run
whose solver status is only feasible. Optimality in that case follows
from the separately proved formula, not the solver status.
The remaining mixed records are positive witnesses; their solver
optimality claims are not independently certified. Records may repeat
an assembly, and the search is not exhaustive. The exact family
formula and these mixed-assembly checks do not prove or disprove
the original conjecture.

\section{Vertex triangles and symmetric bridge obstructions}
\label{sec:round23}
A remaining inductive question is whether replacing a cubic vertex
by a triangle can be the first operation to violate the half-bound.
The exact identity below appeared in the earlier computational
research record; here we use it to isolate a necessary obstruction
and connect it to the symmetric bridge construction.

For a cubic graph $G$ and vertex $v$, let $G^\triangle_v$ replace
$v$ by a triangle, with its three new vertices attached bijectively
to the three former neighbours of $v$. Define
\[
\begin{split}
 B_v&=\max\{|S|:S\text{ good in }G,\ v\notin S\},\\
 C_v&=\max\{|S|:S\text{ good in }G,\ S\cap N_G[v]=\varnothing\},\\
 D_v&=\max_{e\ni v}\max\{|S|:S\text{ good in }G-e,\ v\in S\}.
\end{split}
\]
An infeasible maximum is $-\infty$.
\begin{lemma}\label{lem:vertex-triangle-exact}
With this notation,
\[
 c_{\rm ind}(G^\triangle_v)=\max\{B_v,C_v+3,D_v+1\}.
\]
\end{lemma}
\begin{proof}
A good set meets the new triangle in zero, two, or three vertices;
a single selected triangle vertex cannot have two selected neighbours.
The zero-vertex case gives exactly $B_v$. In the three-vertex case,
all outside neighbours of the triangle must be unselected, giving
$C_v+3$. In the two-vertex case, both corresponding outside neighbours
must be selected. Contract the selected triangle edge back to $v$
and delete the original edge toward the unselected triangle vertex's
outside neighbour. The result is good in $G-e$ and contains $v$.
Conversely, every such good set in $G-e$ lifts using the two triangle
vertices corresponding to the retained edges at $v$, gaining one
vertex. The outside neighbour on the deleted edge may be selected,
which is why $G-e$ cannot be replaced by $G$ in the definition of
$D_v$.
\end{proof}

\begin{theorem}\label{thm:vertex-triangle-obstruction}
Let $|V(G)|=2h$ and suppose $c_{\rm ind}(G)\geq h$.
The graph $G^\triangle_v$ violates \eqref{eq:conj} if and only if
\[
 B_v=h,\qquad D_v\leq h-1,\qquad C_v\leq h-3.
\]
In that case $c_{\rm ind}(G)=h$ and
$c_{\rm ind}(G^\triangle_v)=h$; no maximum good set of $G$ contains
$v$. Moreover, for every edge $e$ incident with $v$, subdividing $e$
once to form a rooted graph $Q_e$ and joining two copies at their
roots gives
\[
 |V(D(Q_e))|=4h+2,\qquad c_{\rm ind}(D(Q_e))=2h.
\]
Thus every such failed triangle expansion also produces a symmetric
bridged counterexample to the original conjecture.
\end{theorem}
\begin{proof}
Every good set of $G$ containing $v$ remains good after deletion of
its unique incident edge toward an unselected neighbour. Its size
is therefore at most $D_v$. Consequently
$c_{\rm ind}(G)\leq\max\{B_v,D_v\}$.
If the expansion violates the half-bound, its order is $2h+2$ and
Lemma~\ref{lem:vertex-triangle-exact} gives
$B_v\leq h$, $D_v\leq h-1$, and $C_v\leq h-3$.
The assumed lower bound on $G$ forces $B_v=h$ and
$c_{\rm ind}(G)=h$. The converse and the exact expanded value follow
directly from the same identity. The bound on $D_v$ excludes $v$
from every maximum good set.

Fix an incident edge $e=vu$. In $G-e$, a good set omitting $v$ has
maximum size $B_v=h$, while one containing $v$ has size at most
$D_v\leq h-1$. Hence $c_{\rm ind}(G-e)=h$. A good set of $G$
containing both endpoints of $e$ has size at most $h-1$, since it
contains $v$. After subdivision, the root-out optimum is therefore
$B=h$, and the root-in optimum is $A\leq h$. The exact doubling
identity of Lemma~\ref{lem:double} yields
$\max\{2B,A+B\}=2h$, on $4h+2$ vertices.
\end{proof}

\begin{corollary}\label{cor:triangle-covered}
If $c_{\rm ind}(G)>|V(G)|/2$, every single vertex-to-triangle
expansion satisfies \eqref{eq:conj}. The same is true at a vertex
belonging to a good set of size at least $|V(G)|/2$. In particular,
if $V(G)$ partitions into two good sets, every single such expansion
satisfies the bound.
\end{corollary}
\begin{proof}
For a strict half-bound witness, keep it if $v$ is absent and lift it
with a gain of one if $v$ is present. A half-sized witness containing
$v$ also lifts with a gain of one. In a partition into two good sets,
every vertex has two neighbours in its own part and one in the other
part. Counting crossing edges shows that the parts have equal size.
Every vertex thus belongs to a half-sized witness.
\end{proof}

\paragraph{Verification and a failed shortcut.}
Independent enumeration checks Lemma~\ref{lem:vertex-triangle-exact}
for all 246 vertex expansions of the 27 connected cubic graphs
through order ten. Separately, complementary good-set partitions
were found and directly verified for 100 previously certified sharp
records: five cubic records and 95 symmetric joins of rooted records.
Their orders range from 24 to 58; distinct records need not be
nonisomorphic. The verifier explicitly expands every vertex and
checks the lifted set, giving 3,932 verified half-bound witnesses.
This does not assert that every sharp cubic graph admits such a
partition, or that repeated expansions preserve it.

One possible shortcut would assert that no cubic vertex can have
all three incident edges increase $c_{\rm ind}$ when deleted. This
assertion is false. The fourteen-vertex graph with graph6 code
\texttt{M???FBOq@KF?D\_B\_?} has $c_{\rm ind}=8$, whereas deleting
any of $\{2,7\}$, $\{2,10\}$, and $\{2,11\}$ raises the value to
10. Exhaustive vertex-subset enumeration independently checks these
values and retains explicit witnesses. The graph satisfies the
original half-bound, since $8>14/2$; it refutes only that auxiliary
shortcut. It does not settle the more restrictive sharp-graph case.

The certificate generator is
\texttt{exploratory/triangle\_obstruction\_certificates.py}, and the
independent verifier is \texttt{verify\_round23.py}. Data are retained
in \texttt{round23\_data}. Theorem~\ref{thm:vertex-triangle-obstruction}
is a conditional reduction, not a proof that its obstruction is
impossible. The original conjecture remains unresolved.

\subsection{Direct searches for the triangle obstruction}
\label{sec:round24}
We next applied the preceding criterion to every cubic parent in
the two radius-four edge neighbourhoods of Section~\ref{sec:round20}.
For each parent, a good set larger than half certifies every single
triangle expansion. Otherwise, half-sized good sets whose union
covers all parent vertices suffice: expand a vertex using a witness
that contains it. The implementation tests previously retained
witnesses and, when necessary, uses exhaustive vertex-assignment
search to find additional ones. If these criteria fail, it directly
checks the expanded graph; no such fallback was needed in the
completed runs.

All single-vertex expansions in both regions satisfy the half-bound.
The 24-vertex region has 1,974,349 reconnection records and hence
47,384,376 indexed expansions. The 26-vertex region has 2,897,883
records and 75,344,958 indexed expansions. In total this is
122,729,334 indexed expansions. These counts include duplicate
constructions and must not be interpreted as numbers of distinct
labelled or nonisomorphic graphs.

The calculation retained 5,656 additional strict parent witnesses
and 76 additional half-bound witnesses. The independent standard-library
verifier checks all 5,732 directly and constructs every triangle
lift justified by each witness: all parent vertices for a strict
witness, and its selected vertices for a half-bound witness. This
gives 142,886 explicit verified lifts. The large reconnection
regions are not independently regenerated; completeness relies on
the C++ enumeration, whose generation scheme was separately checked
against complete small labelled graph lists in Section~\ref{sec:round20}.
The implementation is \texttt{exploratory/triangle\_edge\_ball.cpp}.

A broader constraint-generation search directly targets a cubic
counterexample with one prescribed triangle. Its parents have orders
24 or 26, fixed neighbours at the expanded vertex, and at most six
removed edges relative to a relabelled sharp seed. Each discovered
half-sized good set in the expanded graph becomes a forbidden-set
constraint for subsequent candidates. The two runs produced eight
and six candidates, accumulating 1,821 and 1,907 constraints, then
stopped at wall-time limits. Independent cycle-space enumeration
checks every candidate's reported maximum and number of half-bound
witnesses; it also verifies that each parent satisfies the original
bound. The timeouts exclude no additional region. The source is
\texttt{exploratory/triangle\_counterexample\_synthesis.py}.

Finally, all 21 edge deletions of the fourteen-vertex critical-star
example were joined to the 27 retained two-port profiles, in both
orientations. All 1,134 boundary-table checks have doubled surplus
at least two. Thus this particular attempt to turn the auxiliary
critical-star example into a sharp graph or an original counterexample
also fails. It is a finite fragment calculation, not a theorem about
arbitrary partners.

The original conjecture remains unresolved; see \texttt{round24\_verification.json}.

\section{Repair barriers with only a total edge bound}
The pseudoforest restriction is not responsible for all the barriers in
Section~13. Allow arbitrary odd spanning factors as intermediate states,
subject only to $|E(M)|\leq n$. This is the relevant edge bound because
an isolated-edge-free odd factor with at most $n$ edges gives a good set
of size at least $n/2$. Define $h^{\rm edge}_G(M)$ as the minimum of
$\max_i j(M_i)-j(M)$ over such simple-cycle exchange sequences ending
at a factor with smaller $j$. There is no restriction on the cycle
ranks of their components.

\begin{theorem}\label{thm:edge-budget-barrier}
For the graphs $B_k$ and initial factors $M_k$ of Section~13,
\[
 \left\lceil\frac{k-1}{4}\right\rceil
 \leq h^{\rm edge}_{B_k}(M_k)\leq k-1 \qquad(k\geq3).
\]
In particular, allowing components with arbitrarily many cycles does
not permit repairs with any fixed bound on the temporary increase in
isolated edges. Nevertheless, $\cind(B_k)=|V(B_k)|/2$.
\end{theorem}
\begin{proof}
Consider any restriction $N$ of an odd factor to a cap before either
displayed $k$-cycle has been exchanged. Its root degree is zero or two,
and its other twelve degrees are odd. Retain the notation
$w(N)=j_P(N)+\mathbf1_{\{\deg_N(r)=0\}}$, but impose no
pseudoforest condition. We claim
\begin{equation}\label{eq:cap-budget}
 |E(N)|\geq13-4w(N).
\end{equation}
If $w(N)=0$, the root has degree two and there is no isolated edge.
The nonroot degree-one vertices $L$ induce a good set in $P$:
edges outside $N$ give two neighbours in $L$ to each such vertex,
and an edge of $N$ between two vertices of $L$ would be isolated.
The cap bound gives $|L|\leq6$, and the degree sum gives
$|E(N)|=19-|L|\geq13$. For $w(N)\geq1$, write $d$ for the number of nonroot degree-three
vertices and $\beta\geq0$ for the total cycle rank of $N$.
If the root has degree two, then $m=|E(N)|=7+d$, so $N$ has
$13-m+\beta=6-d+\beta$ components. Apart from the root component,
every component that is not an isolated edge contains a degree-three
vertex. Hence $j_P(N)\geq5-2d+\beta$, giving $2m+w(N)\geq19$.
If the root is isolated, then $m=6+d$ and the nonroot subgraph has
$6-d+\beta$ components. At most $d$ of these are not isolated
edges. Thus $j_P(N)\geq6-2d+\beta$ and again $2m+w(N)\geq19$.
Consequently $m\geq(19-w(N))/2\geq13-4w(N)$ for $w(N)\geq1$.

Every bridge is present in every odd factor: summing odd degrees on
either side of a bridge in a cubic graph forces its inclusion.
Moreover, a simple cycle cannot cross a bridge. Until the first
displayed-cycle exchange, both displayed cycles therefore remain
absent and the central bridge remains an isolated edge. Writing $N_i$
for the $2k-2$ cap restrictions and $W=\sum_iw(N_i)$, we have
\[
 j=1+W,\qquad
 |E(M)|=(2k-1)+\sum_i|E(N_i)|\geq n_k-1-4W.
\]
Some displayed cycle must eventually be exchanged to eliminate the
central isolated edge. Its first exchange adds exactly $k$ edges.
The edge bound immediately after it requires
$n_k-1-4W+k\leq n_k$, whence $4W\geq k-1$.
This proves the lower bound on the peak excess.

The explicit $k$-exchange sequence in Section~13 has edge counts
$n_k-1,n_k-2,\ldots,n_k-k,n_k$ and defect counts
$1,2,\ldots,k,0$. It obeys the total edge bound, proving the
upper bound. The equality for $\cind(B_k)$ was proved there.
\end{proof}

An independent finite check confirms the local bound. Exhausting its $2^{19}$
edge subsets leaves precisely 128 states with the required degree
parities. Their minimum edge counts at each cost are
\[
\begin{array}{c|rrrrrrr}
 w&0&1&2&3&4&5&7\\\hline
 \min |E(N)|&13&9&9&8&8&7&6
\end{array}
\]
The theorem is independent of this enumeration. The table also shows
that the coefficient four in the local bound cannot be reduced: a
state with $w=1$ and nine edges attains equality.
The standard-library program \texttt{verify\_round25.py} checks the
table directly, the root-avoiding cap maximum, and all edges and
degrees in the explicit repairs for $3\leq k\leq40$.

This rules out a bounded-defect cycle-exchange strategy even after
removing the componentwise pseudoforest constraint. It does not
exclude strategies allowing unbounded temporary defects, and it
does not settle the existence of an isolated-edge-free odd factor
with at most $n$ edges in an arbitrary cubic graph. The original
conjecture remains unresolved.

\section{Arbitrary assemblies of four-port fragments}
The following extension of the Hall argument in Section~15 rules out
another direct construction of counterexamples. A four-port fragment
$R$ has four distinct degree-two vertices, called ports; every other
vertex has degree three. Its order $q$ is even. The fragment need not
be connected. Call a set \emph{closed} if it induces a $2$-regular
subgraph of $R$, so that every selected port already has two selected
neighbours within $R$.

\begin{theorem}[A sufficient condition for arbitrary four-port assemblies]
\label{thm:four-port-assembly}
Suppose each member $R_i$ of a collection of disjoint four-port
fragments, of orders $q_i$, satisfies at least one of these conditions:
\begin{enumerate}
\item There is a closed set of size at least $q_i/2$ avoiding all ports.
\item For each of the six pairs of ports, there is a closed set of size
at least $q_i/2$ whose selected ports are exactly that pair.
\end{enumerate}
Pair all ports by edges between distinct fragments, obtaining a simple
cubic graph $G$. Then $\cind(G)\geq |V(G)|/2$, for any number of
fragments and any such pairing.
\end{theorem}
\begin{proof}
Use the first condition wherever available, and call the remaining
fragments active. Make a bipartite graph with two demand vertices for
each active fragment on the left, and one resource vertex for each
joining edge on the right. Both demands of a fragment are adjacent
to its four joining edges. For any set of demands, let $s$ be the
number of represented fragments. Their $4s$ port incidences belong
to at least $2s$ distinct joining edges, because each edge has at most
two ends in these fragments. The number of demands is at most $2s$.
Hall's condition therefore holds, and a matching assigns two distinct
joining edges to every active fragment, with no edge assigned twice.

In an active fragment select the closed set for the two ports incident
with its assigned edges. In every other fragment select the
port-avoiding closed set. No joining edge has both endpoints selected:
that would require it to be assigned to both endpoint fragments.
Thus the union of these closed sets is good in $G$ and has size at
least $\sum_iq_i/2=|V(G)|/2$.
\end{proof}

We tested fragments obtained by deleting adjacent vertices $u,v$ with
no common neighbour from 56 retained sharp cubic examples: the Franklin
graph, the five examples of orders $24,26,36,38,42$, and the symmetric
doubles of the fifty retained thirteen-vertex caps from Section~18.
These are labelled records, not a claimed nonisomorphic census.
There are 2,214 eligible edge-deletion records. Of these, 1,643 satisfy
the first condition and all remaining 571 satisfy the second.
Consequently \emph{every} assembly made solely from this library
satisfies the conjecture, allowing arbitrary repetition, orientation,
and number of pieces. This conclusion goes beyond a finite search
over pairwise joins, but does not cover arbitrary four-port fragments.

For completeness, the exact profile generator enumerates Eulerian
edge sets of the original cubic graph, removes $u,v$ from their
supports, and checks induced degrees in the fragment. Each port is
unselected (state 0), selected with internal degree two (state 1), or
selected with internal degree one (state 2). The resulting 81-entry
table records the maximum doubled surplus $2|S|-q$ at each state.
This enumeration is exhaustive: a feasible partial cycle system has
zero, two, or four state-2 ports. With two, close its ends through
$u$, $v$, or the path $uv$, according to their locations. With four,
join the two ports formerly adjacent to $u$ through $u$ and the other
two through $v$. With zero, no added edge is needed. In each case
the resulting edge set is Eulerian in the original graph. It need
not be induced there; inducedness is checked only after deletion.

The independent verifier checks all 44,154 feasible-state witnesses
directly and verifies that every eligible deletion of each saved seed
is included. Only positive witnesses are used for
Theorem~\ref{thm:four-port-assembly}; large-profile maxima and
infeasibility claims are not needed for this certificate. As a separate
generator check, all 252 eligible profiles from connected cubic graphs
through order ten agree with exhaustive vertex-subset enumeration.
The verifier also constructs 41 complete assemblies and checks their
good sets directly, including an assembly using every retained piece
with 54,996 vertices. These checks test the construction; its general
validity follows from the Hall proof.
The sources are \texttt{exploratory/four\_pole\_large.cpp} and
\texttt{verify\_round26.py}; results are saved in
\texttt{round26\_verification.json}. The original conjecture remains
unresolved outside the proved classes.

\section{A four-port contraction criterion}
The preceding assembly theorem can be complemented by a local
reduction that applies even when the rest of the graph is arbitrary.
Use the boundary states $0,1,2$ of Section~20, and write $f_s$ for
the maximum doubled surplus $2|S|-q$ in state $s$ of a four-port
fragment of order $q$, with value $-\infty$ for an infeasible state.
The following conditions can equally be certified by individual
witnesses; exact maxima are not needed.

\begin{theorem}[Replacing a fragment by two adjacent vertices]
\label{thm:four-port-contraction}
Let $R$ be an induced four-port fragment in a simple cubic graph $G$,
with ports ordered $p_1,p_2,p_3,p_4$. Suppose
\begin{enumerate}
\item $f_{0000}\geq-2$;
\item $f_{2200}\geq0$ and $f_{0022}\geq0$;
\item for every $i\in\{1,2\}$ and $j\in\{3,4\}$, some state with
entries $s_i=s_j=2$ and the other two entries in $\{0,1\}$ has
$f_s\geq2$.
\end{enumerate}
Delete $R$, add adjacent vertices $u,v$, and reattach the former
external edges at $p_1,p_2$ to $u$ and those at $p_3,p_4$ to $v$.
Assume the resulting graph $G_0$ is simple and cubic. Then
\[
 \cind(G)\geq\cind(G_0)+\frac{q-2}{2}.
\]
Consequently a minimum-order counterexample to the original
conjecture contains no such fragment of order $q>2$ for which the
specified replacement gives a simple cubic graph.
\end{theorem}
\begin{proof}
Fix a good set $S_0$ in $G_0$. If neither $u$ nor $v$ belongs to
$S_0$, replace them by a state-$0000$ set in $R$. It has at least
$(q-2)/2$ vertices and has no selected port, so it has no edge to
the retained selected vertices outside $R$.

If only $u$ is selected, both its external neighbours must be
selected. Use a state-$2200$ set, with at least $q/2$ vertices.
Its two selected ports have internal degree one and receive their
second selected neighbour outside $R$. The other ports are
unselected. All selected exterior vertices retain exactly the
same number of selected neighbours. Replacing the one selected
vertex $u$ therefore gains at least $(q-2)/2$ vertices. The case
where only $v$ is selected uses state $0022$.

If both $u,v$ are selected, each has exactly one selected external
neighbour. Let their corresponding ports be $p_i,p_j$, one in each
pair. Use the state from the third condition. Its open ports attach
to the selected exterior neighbours. Either of the other ports may
be selected with internal degree two, since its exterior neighbour
is unselected. Again all selected degrees are exactly two. This
replaces two selected vertices by at least $(q+2)/2$, giving the
same gain. Maximizing over $S_0$ proves the inequality. Since
$|V(G)|=|V(G_0)|+q-2$, it transfers the conjectured half-bound
from $G_0$ to $G$ and proves the minimal-counterexample assertion.
\end{proof}

With the port pairing inherited from the deleted adjacent vertices,
1,945 of the 2,214 fragments in Section~20 meet this criterion.
The other 269 are not covered by this contraction test. These counts
do not assert that any cubic graph contains a qualifying fragment,
nor that its contraction is always simple. Directly replacing an
edge and its endpoints in $K_{3,3}$ by each certified fragment checks
19,450 lifts, covering all good sets of that six-vertex exterior
model. The general statement follows from the preceding case proof.

We also broadened the direct counterexample search to all eligible
adjacent-vertex deletions from the 621 connected simple cubic graphs
of orders at most fourteen. This produces 9,492 indexed four-port
fragments, whose orders range from four to twelve. Their tables have
7,144 distinct labelled value vectors and 2,054 vectors up to port
permutation. Retaining 1,239 vectors suffices under componentwise
domination and port permutation. Every discarded vector dominates
some retained vector in a suitable port order. Thus a nonnegative
joining witness for the retained vectors transfers to the discarded
ones: replacing either profile by a componentwise larger profile
cannot decrease the best joined surplus.

All $24\binom{1240}{2}=18,436,320$ pairings of retained profiles
with repetition have a nonnegative joining witness. Hence every
two-fragment join from this entire library, in every port orientation,
satisfies the original half-bound. The resulting orders are at most
24. These are indexed constructions, with possible repeated graphs;
they are not a census of cubic graphs on those orders. This finite
result does not imply that arbitrary larger assemblies from this
broader library satisfy the bound.

The verifier \texttt{verify\_round27.py} checks all 210,400 recorded
local witnesses by induced degrees, every discarded-profile domination,
and all retained pairings. Its pair checker uses compatible-state
bitsets grouped by score, independently of the exploratory C++
enumeration of the 256 compatible state pairs for each orientation.
Large-profile upper bounds are not needed for these positive checks.
Data and results are in \texttt{round27\_data/} and
\texttt{round27\_verification.json}. The original conjecture remains
unresolved.

\section{Assemblies using open ports and a two-factor}
The closed-port condition of Section~20 is not the only way to
assemble fragments. Keep the same states and doubled-surplus
profiles. For a pair $I$ of ports, let $s(I)$ be the state with
entries two at $I$ and zero at the other ports. Write $e_j$ for
the state vector with entry one only at port $j$.

\begin{theorem}[Open-port assembly criterion]\label{thm:open-port-assembly}
Suppose every fragment $R$ in a collection of four-port fragments
satisfies the following condition for every pair $I$ of its ports:
\[
 f_{s(I)}\geq0
 \quad\text{or}\quad
 f_{s(I)+e_j}\geq0\ \text{for both }j\notin I.
\]
Any simple cubic graph obtained by pairing all ports between
distinct fragments satisfies $\cind(G)\geq |V(G)|/2$.
\end{theorem}
\begin{proof}
Contract each fragment to a vertex, retaining the joining edges.
The resulting loopless multigraph $Q$ is $4$-regular. It has a
$2$-factor $F$: orient an Euler tour in each component, giving
indegree and outdegree two, and take a perfect matching in the
associated bipartite graph with one outgoing and one incoming copy
of each vertex. This bipartite graph is $2$-regular, so Hall's theorem
provides the matching. Its directed edges give the required $2$-factor.

In each fragment the two ports incident with $F$ will have state two.
If state $s(I)$ has nonnegative surplus, use it. Call the remaining
fragments active. Each active fragment can instead use either of
its two non-$F$ ports in state one, leaving the other in state zero.
Assign one incident non-$F$ edge to each active fragment, without
assigning any edge twice. Hall's theorem again applies: any $s$
active fragments have $2s$ incidences with non-$F$ edges, belonging
to at least $s$ distinct edges.

Use the state-one port on the assigned edge. A joining edge in $F$
has state two at both ends; it completes each selected endpoint's
internal degree one to degree two. A joining edge outside $F$ has
at most one selected endpoint, since it is assigned at most once.
It therefore adds no selected neighbour to a state-one port, whose
internal degree is already two. The union of the local sets is good,
and every fragment contributes at least half its vertices.
\end{proof}

Of the 9,492 fragments in Section~21, 5,552 satisfy this condition.
Thus arbitrary assemblies using only those fragments satisfy the
half-bound. The independent verifier checks their local witnesses
and constructs 21 whole-graph examples using Euler orientations and
two bipartite matchings. One uses every qualifying fragment and has
65,380 vertices. Its selected set is checked by actual induced
degrees. This is a second sufficient assembly criterion; arbitrary
mixtures with fragments satisfying only Theorem~20.1 are not claimed
to be covered.

\subsection{Direct cyclic-assembly searches}
We also searched assemblies whose incidence multigraph is a cycle
with two parallel edges between consecutive fragments. Each fragment
has two left and two right ports. If $a,b\in\{0,1,2\}^2$, define
the $9$-state transfer matrix
\[
 M_R(a,b)=\max_{t\sim a} f_R(t,b),
\]
where compatibility is coordinatewise and permits exactly the pairs
$(0,0),(0,1),(1,0),(2,2)$. With max-plus matrix multiplication,
the optimum doubled surplus of a fixed cyclic assembly is
\[
 \max_a (M_{R_1}M_{R_2}\cdots M_{R_k})(a,a).
\]
This is exact because a compatible state choice determines disjoint
local solutions, and the boundary conditions make their union good.

The retained 1,239 profiles and all 24 port permutations give 29,736
indexed choices per position. Coordinate descent scans all choices
at one position while keeping the others fixed. The recorded search
uses $k\in\{3,4,5,6,8,12,20,40\}$, three restarts per length and
five sweeps per restart: 43,711,920 coordinate evaluations and
1,470 retained assembly records. This is a heuristic search over
words of fragments, not exhaustive enumeration of all such words.
Independent integer dynamic programming reconstructs and verifies
every recorded value and a whole-graph good set. Their orders range
from 12 to 470, and none violates the half-bound.

All eight zero-surplus words are Franklin relabellings, verified by
explicit isomorphisms. CP-SAT separately excludes larger good sets;
complementary good partitions are checked by induced degrees. Data
are in \texttt{round28\_data/}; independent checks are implemented in
\texttt{verify\_round28.py} and \texttt{verify\_round28\_assembly.py}.
The original conjecture remains unresolved.

\subsection{Mixing the two assembly criteria}
The two criteria can be combined when the joining graph supplies
a suitable collection of cycles. This is a sufficient condition,
not an assertion that all mixtures are covered.

\begin{theorem}[Mixed assembly criterion]\label{thm:mixed-assembly}
Let $Q$ be the loopless $4$-regular multigraph of joining edges of
a four-port assembly. Suppose $Q$ has a subgraph $F$ in which each
vertex has degree zero or two. At every degree-zero vertex, require
the closed-port condition of Theorem~20.1. At every degree-two
vertex, require the open-port condition of Theorem~22.1. Then
the assembled simple cubic graph $G$ satisfies
$\cind(G)\geq |V(G)|/2$.
\end{theorem}
\begin{proof}
For a degree-zero fragment, choose its port-avoiding solution if
available, creating no demand. Otherwise it can select any pair
of ports and demands two distinct joining edges. For a degree-two
fragment, give state two to the ports on $F$. It either has a
nonnegative-surplus solution avoiding both remaining ports, or can
use either remaining port in state one and demands one non-$F$ edge.

All demands are assigned to edges of $Q-F$, with no edge assigned
twice. Indeed, at a degree-zero vertex $Q-F$ has degree four and
there are at most two demands; at a degree-two vertex it has degree
two and there is at most one demand. For any set of demand slots,
let $W$ be the represented vertices. The number of slots is at most
$\tfrac12\sum_{v\in W}d_{Q-F}(v)$, whereas the number of incident
resource edges is at least this quantity. Hall's theorem supplies
the assignment. Use exactly the assigned ports in state one.

Edges of $F$ have state two at both ends and complete the selected
degrees to two. Every other joining edge has at most one endpoint
in state one and neither endpoint in state two. Thus all selected
degrees are two, and the local nonnegative surpluses sum to a
nonnegative surplus for $G$.
\end{proof}

Taking $F$ empty recovers the closed-port theorem; taking it spanning
recovers the open-port theorem. In particular, for a $K_5$ joining
graph, any mixture with three or four designated open-type fragments
is covered: take a cycle through exactly those vertices. With only
one or two designated open-type fragments, this choice of $F$ is
impossible. This failure of the sufficient condition is not a
counterexample to the original conjecture. The verifier constructs
and degree-checks 1,000 mixed examples of orders 44 through 60,
using fragments that satisfy only their designated criterion.

A separate counterexample-guided search fixes the $K_5$ joining
graph and allows ten closed-type and ten open-type profiles, each
in all 24 port orders. It fixes one block to each type, to require
a genuine mixture. Each nonnegative whole-graph witness supplies
an exclusion constraint on the profile choices. After 120.7 seconds
the run reached its time limit with 41 such witnesses and 78 solver
queries. The witnesses have orders 36 through 60 and doubled
surplus at least twelve. The time limit does not prove the remaining
restricted domain infeasible. All saved witnesses and their boundary
compatibilities are checked directly by \texttt{verify\_round29.py}.

Finally, a focused exhaustive check uses only the four-cycle fragment
and the eight-vertex fragment obtained from graph6
\texttt{I?BeeOwM?} by deleting vertices 2 and 5 (zero-based labels).
The latter has $f_{0000}=-8$ and all six closed-pair values at least
zero. Each piece has three distinct boundary profiles under its 24
port permutations. For one or two four-cycle pieces and all remaining
pieces of the second type on $K_5$, there are $2\cdot3^5=486$
indexed profile assignments. Every one has a directly checked good
set with doubled surplus at least ten. The graphs have orders 32 or
36. This exhausts these two selected fragment types on $K_5$, not
the wider mixed library. The source is
\texttt{exploratory/mixed\_k5\_exhaust.py}; the original conjecture
remains unresolved.

\section{Two-edge extensions and the order-22 census}
Let $e=uv$ and $f=wz$ be disjoint edges of a simple cubic graph $G$.
Let $G_{e,f}^{+}$ be obtained by subdividing $e,f$ with new vertices
$x,y$, respectively, and adding $xy$. This increases the order by two
and preserves simplicity and cubicity. Unlike vertex-triangle expansion,
the operation need not introduce a triangle.

Define $B=\cind(G-e-f)$. Let $T_e$ be the largest good-set size in
$G-f$ subject to containing both $u,v$, and define $T_f$ symmetrically
in $G-e$. Finally, let $P$ be the largest size of a vertex set $S$ in
$G-e-f$ such that exactly one endpoint of each of $e,f$ belongs to
$S$, these two selected endpoints have induced degree one, and every
other vertex of $S$ has induced degree two. Thus $S$ induces a path
between the two selected endpoints, together with possibly some cycles.
Infeasible constrained maxima have value $-\infty$.

\begin{proposition}[Exact two-edge extension identity]
\label{prop:two-edge-extension}
\[
 \cind(G_{e,f}^{+})=\max\{B,T_e+1,T_f+1,P+2\}.
\]
\end{proposition}
\begin{proof}
Separate good sets according to their intersection with $\{x,y\}$.
If neither is selected, their restriction to the original vertices is
exactly a good set of $G-e-f$. If only $x$ is selected, both $u,v$
must be selected; restoring the edge $uv$ in place of the selected
path $uxv$ gives a good set of $G-f$. This correspondence is reversible
and changes the size by one. The case where only $y$ is selected is
symmetric. If both are selected, each has exactly one selected old
neighbour. Removing $x,y$ leaves precisely a set counted by $P$.
Conversely, adding $x,y$ and their joining edge closes its path into
an induced cycle, leaving its other cycle components unchanged. This
last correspondence changes the size by two.
\end{proof}

For $|V(G)|=2h$, this construction would violate the original
conjecture exactly when
\[
 B\leq h,\qquad T_e,T_f\leq h-1,\qquad P\leq h-2.
\]
These are exact conditions for this construction, not their exclusion.
We tested every disjoint edge pair in 101 retained sharp parent records:
the Franklin graph and the 100 records used in Section~18. All 182,679
extensions have directly verified half-bound witnesses. Their orders
range from 14 to 60; repeated constructions are possible. Independently,
the displayed identity agrees with exhaustive vertex-subset enumeration
in all 1,674 eligible edge-pair cases from connected cubic graphs through
order ten. Data are in \texttt{round30\_data/}; the verifier is
\texttt{verify\_round30\_extensions.py}.

\subsection{Complete connected census at order twenty-two}
The unrestricted computation at order 22 uses nauty 2.9.3 with command
\texttt{geng -c -d3 -D3 -X0 -x3200 22 r/64}, for all residues
$r=0,\ldots,63$. The 64 disjoint shards contain 7,319,447 connected
simple cubic graphs up to isomorphism, agreeing with the published
House of Graphs count~\cite{HoGCubic}. For every graph we retain a
vertex mask whose selected vertices induce a $2$-regular graph of
order at least twelve. The masks and their graph6 inputs are stored
in compressed per-shard certificate files.

A separate C++ checker imports no search routines. It decodes each
graph, checks all degrees and connectivity, and directly counts the
selected neighbours of every selected vertex. Every certificate passes.
The generator's count, the search count and the independent checker
count agree for every shard. Enumeration completeness relies on
\texttt{geng}'s isomorph-free generation and documented residue splitting;
the degree check independently establishes the positive bound for every
retained graph. This is a computational census, not a general proof.

The bound twelve is attained. The graph with graph6 string
\begin{center}\texttt{U?????????[?s?q?Co@o?DO?K\_?Co?F??@W??M??}\end{center}
has a good set
$\{0,2,3,6,7,9,11,12,14,16,18,20\}$, using zero-based labels.
Independent enumeration of all $2^{22}=4,194,304$ vertex subsets gives
\[
\begin{array}{c|rrrrrr}
 |S|&0&4&6&8&10&12\\\hline
 \#\text{ good sets}&1&6&18&24&66&84
\end{array}
\]
and no larger good set. Thus the minimum of $\cind(G)$ over connected
simple cubic graphs of order 22 is exactly twelve. Together with the
previous smaller-order checks, any counterexample to the original
conjecture must have at least 24 vertices: a disconnected counterexample
would have a connected component that is itself a counterexample.
The programs \texttt{verify\_round30\_census.py} and
\texttt{verify\_round30\_minimum.py} record the checks and certificate
hashes. No proof or counterexample for arbitrary orders has been found.

\section{A two-edge extension can lose arbitrarily many vertices}
The exact identity in Section~23 does not imply an additive induction
step. In fact the optimum can decrease without any uniform bound under
a single two-edge extension. We give an explicit infinite family and
explain why it still satisfies the original conjecture.

\subsection{An attachment that preserves the unweighted optimum}
Let $Q$ be the triangular prism with triangles $abc$ and $a'b'c'$ and
matching edges $aa',bb',cc'$, with the edge $ab$ subdivided by a root $z$.
The root has degree two and the other six vertices have degree three.

\begin{lemma}\label{lem:neutral-cap}
The largest good sets in $Q$ containing and avoiding $z$ both have size five.
\end{lemma}
\begin{proof}
The cycles $azbb'a'a$ and $acbb'a'a$ are induced five-cycles, containing
and avoiding $z$, respectively. A six-vertex good set would contain at
least five of the six degree-three vertices. Each of those vertices
would need a neighbour outside the set. The unique outside vertex has
degree at most three, a contradiction. The full seven-vertex set is
not good either. Thus five is an upper bound as well as an attained value.
\end{proof}

For a subcubic graph $H$, attach a separate copy of $Q$ at each missing
incidence, joining its root to the deficient vertex. Write $\mathcal N(H)$
for the resulting cubic graph and $D=3|V(H)|-2|E(H)|$.

\begin{proposition}\label{prop:neutral-completion}
For every subcubic graph $H$,
\[
 |V(\mathcal N(H))|=|V(H)|+7D,\qquad
 \cind(\mathcal N(H))=5D+\cind(H).
\]
\end{proposition}
\begin{proof}
Each new joining edge is a bridge, so no good set uses it. The
restriction of a good set to $H$ is good, and each attached copy
contributes at most five vertices by Lemma~\ref{lem:neutral-cap}.
Conversely, any good set in $H$ extends by choosing a root-avoiding
five-cycle in every copy. This is valid whether the corresponding
vertex of $H$ is selected or not. Both bounds are attained together.
\end{proof}

This differs from the branch-counting completion in Section~3:
it preserves the unweighted optimum up to an additive constant.
It is useful for transferring sensitivity examples from subcubic cores
to cubic graphs, but its surplus does not preserve the conjectured
half bound.

\subsection{An exact unbounded loss}
For $k\geq2$, form $G_k$ by taking a cycle
$v_0v_1\cdots v_{2k-1}v_0$ and attaching a copy of $Q$ at each vertex.
Let $G_k^+$ be its two-edge extension at the disjoint edges
$v_0v_1$ and $v_kv_{k+1}$, with new vertices $x,y$ and edge $xy$.

\begin{theorem}\label{thm:unbounded-extension-loss}
These graphs are connected, simple and cubic, and satisfy
\begin{align*}
 |V(G_k)|&=16k,&\cind(G_k)&=12k,\\
 |V(G_k^+)|&=16k+2,&\cind(G_k^+)&=11k+2.
\end{align*}
Consequently $\cind(G_k)-\cind(G_k^+)=k-2$, which is unbounded.
\end{theorem}
\begin{proof}
There are $2k$ attachments, each contributing at most five vertices,
and a root-avoiding optimum is always available in each. Their joining
edges remain bridges after the extension. Thus in both graphs the
optimum is $10k$ plus the optimum of the remaining core.

The original core is an induced cycle on $2k$ vertices. The new core
is a theta graph with three internally disjoint $x$--$y$ paths: the
edge $xy$ and two paths each of length $k+1$. Every cycle in a theta
graph uses two of its three paths. The cycle using both long paths
is not induced, because $xy$ is a chord. Each cycle using $xy$ and
one long path is induced and has $k+2$ vertices. No two cycles are
vertex-disjoint, since both would contain $x,y$. Hence the largest
induced $2$-regular subgraph of the core has exactly $k+2$ vertices.
Adding the attachment contributions gives the two formulas.
\end{proof}

More generally, if a cycle has length $L$ and the two subdivided
edges have cyclic separation $d$, where $2\leq d\leq L-2$, the
extended core has optimum $\max(d,L-d)+2$. This gives a loss of
$\min(d,L-d)-2$ after attaching the same gadgets.

The theorem refutes both the proposed inequality
$\cind(G_{e,f}^+)\geq\cind(G)+1$ and any proposed universal lower
bound $\cind(G_{e,f}^+)\geq\cind(G)-C$ with a fixed constant $C$.
It does not refute the original conjecture: the exact surplus in the
extended family is
\[
 \cind(G_k^+)-|V(G_k^+)|/2=3k+1>0.
\]
An induction based on either of these additive inequalities is therefore
invalid. The boundary quantities in Section~23 retain information that
those inequalities discard; no sufficient general induction step has
been established here.

\subsection{Independent checks and smaller examples}
Exhaustive vertex-subset checks give two smaller failures of additive
monotonicity. Edge endpoints use zero-based graph6 labels.
\begin{center}
\begin{tabular}{@{}lcccc@{}}
\toprule
Parent graph6 & Extended edges & Orders & Before & After\\
\midrule
\texttt{K?ABCN\_U?wF?} & $(0,7),(3,10)$ & $12\to14$ & $9$ & $8$\\
\texttt{O???C@\_sAaSGIGB\_@o?k?} & $(0,7),(8,12)$ & $16\to18$ & $12$ & $10$\\
\bottomrule
\end{tabular}
\end{center}
Here ``before'' and ``after'' are exact maxima, not just witness sizes.
Both extended graphs satisfy the original half bound.
The program \texttt{verify\_round31.py} independently enumerates all
subsets of these four graphs and of the seven-vertex attachment.
It also checks the theta-core formula exhaustively through core order
$L=18$ and constructs whole-graph witnesses for $2\leq k\leq100$.
The infinite-family upper bound follows from the preceding proof,
not from extrapolating these computations. The original conjecture
remains unresolved.

\section{Edge-boundary certificates for bridge joins}
A decrease in the full optimum under an extension does not by itself
threaten the half bound. We now test a sufficient boundary condition
that certifies that bound directly for rooted subdivisions and their
bridge joins. The condition is not proved for arbitrary cubic graphs.

\begin{proposition}\label{prop:edge-half-certificate}
Let $G$ be a connected simple cubic graph of order $2h$, and let $S$
be a good set of size $c\geq h+1$. Suppose that for every edge
$e\in E(G[S])$ there is a good set $T_e$ in $G-e$ with
\[
 |T_e|\geq 2h-c.
\]
For every edge $e$ of $G$, let $Q_e$ be obtained by subdividing $e$
once and taking its new vertex as the root. Then
$\cind(D(Q_e))\geq 2h+1$ for every $e$. Moreover, joining the roots of
any two fragments certified in this way, possibly from different
parent graphs of different orders, satisfies the original half bound.
\end{proposition}
\begin{proof}
If $e\notin E(G[S])$, subdividing it leaves $S$ good with the root
unselected. Hence the root-out optimum $B_e$ is at least
$c\geq h+1$, and $2B_e\geq2h+2>2h+1$.

If $e\in E(G[S])$, subdividing its selected edge gives a root-in
set of size $c+1$. The set $T_e$ gives a root-out set of size at
least $2h-c$. Thus
\[
 A_e+B_e\geq(c+1)+(2h-c)=2h+1.
\]
In either case the doubling formula in Lemma~\ref{lem:double}
certifies the half bound for $D(Q_e)$. The final assertion follows
from bridge symmetrization, Theorem~\ref{thm:bridge-symmetry}:
if each individual symmetric join meets its half bound, so does
any mixed join of the two rooted fragments.
\end{proof}

The strict inequality $c>h$ is essential to this particular
certificate: when $c=h$ and $e\notin E(G[S])$, the available root-out
bound gives only $2B_e\geq2h$, one below the required order $2h+1$.
The proposition does not assert that its boundary hypothesis always
holds, even when a strict good set is known.

\subsection{Exhaustive verification for order-22 parents}
We applied Proposition~\ref{prop:edge-half-certificate} to the saved
good set in every one of the 7,319,447 connected simple cubic graphs
on 22 vertices from Section~23.1. Each saved set has size at least
12. Every required edge-deletion witness was found and independently
checked. There are 33 edges in each parent, hence 241,541,751 indexed
parent-edge records in total. Only edges inside the saved good set
require an additional mask; for every other edge that same good set
already certifies a sufficiently large root-out state.

The producer \texttt{exploratory/edge\_half\_certificates.cpp}
writes the parent graph6 string, its original mask and one deletion
mask for each selected edge, in graph6 edge order. A separate checker,
\texttt{exploratory/check\_edge\_half.cpp}, imports no search code.
It independently decodes each graph, checks simplicity through the
encoding, degree three, connectivity and the induced degree of every
selected vertex. For each deletion mask it removes exactly the
specified edge before checking degree two and the required cardinality.
The Python driver \texttt{verify\_round32.py} also compares every
parent graph and original mask, row for row, against the retained
order-22 census. All 64 shards pass. Enumeration completeness is
inherited from the nauty census; the new witnesses are checked
without trusting the optimization routine. In total, 95,430,602 new
deletion masks are checked, and 146,111,149 parent-edge records reuse
their original mask. An additional direct construction verifies good
sets in 2,048 sampled mixed joins on 46 vertices.

\begin{corollary}[Finite rooted family]\label{cor:rooted23}
Let $Q_1,Q_2$ be connected simple graphs on 23 vertices, each with
one designated vertex of degree two and all other vertices of
degree three. Suppose the two neighbours of each designated root
are nonadjacent. Joining the roots of $Q_1,Q_2$ by an edge produces
a cubic graph of order 46 with an induced $2$-regular subgraph on
at least 23 vertices.
\end{corollary}
\begin{proof}
Suppressing either root produces a connected simple cubic graph
on 22 vertices: its neighbours are distinct and nonadjacent, so
suppression creates neither a loop nor a parallel edge. Each
fragment is therefore an edge subdivision of a parent covered by
the exhaustive check. Apply Proposition~\ref{prop:edge-half-certificate}.
\end{proof}

This corollary is a computationally certified finite-family result.
It includes symmetric and mixed joins, but not all cubic graphs of
order 46. In particular, it excludes roots whose two neighbours
are adjacent, and it gives no unrestricted assertion at higher
orders. No counterexample to the original conjecture was obtained.

\section{A normal form for rooted bridge obstructions}
We now remove the triangle-at-the-root restriction from the finite
result in Section~25. The structural reduction in this section applies
at arbitrary orders; only its final census consequence is finite.

Let $Q$ be a connected simple graph with one root $r$ of degree two
and every other vertex of degree three. Define
\[
 A=\max_{S\text{ good},\ r\in S}|S|,\quad
 B=\max_{S\text{ good},\ r\notin S}|S|,\quad
 C=\cind(Q-N_Q[r]),
\]
using $A=-\infty$ when necessary. Put $d(Q)=\cind(D(Q))=
\max(2B,A+B)$. A \emph{rooted obstruction} means $d(Q)<|V(Q)|$.
This is a counterexample test for the cubic double, not a new
version of the original conjecture.

\subsection{Triangle expansion at the root}
Replace the root $r$, whose neighbours are $a,b$, by a triangle
$r'p q r'$, adding $ap$ and $bq$, and take $r'$ as the new root.
Denote this operation by $E$; it adds two vertices.

\begin{lemma}[Exact root expansion]\label{lem:root-expansion-map}
The profile of $E(Q)$ is
\[
 (A',B',C')=(C+3,\ \max(B,A+1),\ B).
\]
Consequently the profile of $E^2(Q)$ is
\[
 (A'',B'',C'')=
 (B+3,\ \max(B,A+1,C+4),\ \max(B,A+1)).
\]
\end{lemma}
\begin{proof}
A good set containing $r'$ must contain $p,q$ and avoid $a,b$.
The remaining choices are precisely the good sets counted by $C$,
so $A'=C+3$.
If $r'$ is absent, either both $p,q$ are absent or both are present:
exactly one cannot have two selected neighbours. The first case
contributes $B$. In the second, $a,b$ must both be selected;
replacing the selected path $apqb$ by $arb$ gives a root-in good
set of $Q$, smaller by one. Thus $B'=\max(B,A+1)$.
Finally, deleting the closed neighbourhood of $r'$ leaves $Q-r$,
whose optimum is $B$. Applying the same map again proves the
second formula.
\end{proof}

\begin{proposition}[Two expansions preserve the half bound]
\label{prop:root-two-expansions}
If $d(Q)\geq q=|V(Q)|$, then $d(E^2(Q))\geq q+4$.
\end{proposition}
\begin{proof}
If $A>B$, then $d(Q)=A+B$, while
$A''+B''\geq(B+3)+(A+1)=d(Q)+4$.
If $A\leq B$, including $A=-\infty$, then $d(Q)=2B$.
The order $q$ is odd, so $2B\geq q$ implies $2B\geq q+1$.
Since $A''+B''\geq(B+3)+B$, the new double has optimum at least
$q+4$ in this case too. Its required half bound is $q+4$.
\end{proof}

This does not assert that a single root expansion preserves the bound.
For a non-obstructing $Q$ of order $2h+1$, the exact failure criterion is
\[
 E(Q)\text{ is a rooted obstruction}
 \quad\Longleftrightarrow\quad
 B=h+1,\qquad A\leq h,\qquad C\leq h-2.
\]
Indeed, failure forces $B'=\max(B,A+1)\leq h+1$, hence $A\leq h$.
If $B\leq h$ too, the original double would already fail, so $B=h+1$.
Now $A'+B'=C+h+4<2h+3$ is exactly $C\leq h-2$.
Conversely these three conditions make both $2B'$ and $A'+B'$
smaller than $2h+3$. They describe a potential boundary obstruction;
no graph realizing a counterexample profile has been found.

\subsection{The diamond termination case}
Given a rooted graph $R$ with root $b$, add vertices $r,u,v,a$
and edges
\[
 ru,rv,uv,ua,va,ab,
\]
and designate $r$ as the root. Write $L(R)$ for the resulting rooted
graph. Its joining edge $ab$ is a bridge, and its order is $|V(R)|+4$.

\begin{lemma}\label{lem:root-diamond}
If $R$ has profile $(A,B,C)$ and $M=\max(A,B)$, then $L(R)$ has
profile
\[
 (M+3,\ \max(B+3,M),\ M).
\]
Moreover, $d(L(R))\geq d(R)+6$, and
\[
 d(E(L(R)))=2M+8.
\]
Thus if $R$ is not a rooted obstruction, neither $L(R)$ nor
$E(L(R))$ is one.
\end{lemma}
\begin{proof}
No selected cycle can use the bridge $ab$. When $r$ is selected,
the triangle $ruv$ is selected and $a$ is absent, so any maximum
good set in $R$ is available. This gives $M+3$.
When $r$ is absent, either no diamond vertex is selected, contributing
$M$, or the triangle $auv$ is selected, requiring $b$ to be absent
and contributing $B+3$. Removing $N[r]$ leaves $a$ with degree one;
it cannot be selected, leaving optimum $M$ in $R$.

Since $d(R)=M+B$, selecting a root-in optimum in one copy of $L(R)$
and a root-out optimum in the other gives
$d(L(R))\geq(M+3)+(B+3)=d(R)+6$.
The expansion map gives root-in optimum $M+3$ and root-out optimum
$M+4$ for $E(L(R))$, proving its displayed double optimum.
Finally $2M\geq d(R)$; the orders of $L(R)$ and $E(L(R))$ exceed
that of $R$ by four and six, respectively. These inequalities prove
the last assertion.
\end{proof}

\begin{theorem}[Rooted normal form]\label{thm:root-normal-form}
If a rooted obstruction exists, one of minimum order is obtained
from a connected simple cubic graph by one of the following operations:
\begin{enumerate}
\item Subdivide one edge once and use the new vertex as root.
\item Replace an edge $ab$ by $ap,pq,qb,pr,qr$, with fresh vertices
$p,q,r$, and use $r$ as root.
\end{enumerate}
Consequently, if any connected cubic counterexample with a bridge
exists, a counterexample exists among symmetric doubles of these
two kinds of rooted graphs.
\end{theorem}
\begin{proof}
Choose a rooted obstruction $Q$ of minimum order. If its root's
neighbours are nonadjacent, suppress the root. The result is a
connected simple cubic graph, giving the first form.

Otherwise its root $r$ belongs to a triangle $ruv$. Let $a,b$ be
the third neighbours of $u,v$. If $a=b$, this is precisely $L(R)$
for the smaller rooted component $R$ on the other side of the
bridge incident with $a$. Minimality and Lemma~\ref{lem:root-diamond}
give a contradiction.

If $a\ne b$, contract the triangle to a single degree-two root,
obtaining a smaller connected simple rooted graph $Q'$ with
$Q=E(Q')$. If the neighbours of the root of $Q'$ are nonadjacent,
suppress that root to obtain the second stated form.
If they are adjacent, inspect their third neighbours. When those
are distinct, $Q'=E(Q'')$ for a smaller rooted graph $Q''$.
Minimality makes $Q''$ non-obstructing, and
Proposition~\ref{prop:root-two-expansions} contradicts $Q=E^2(Q'')$.
When those third neighbours coincide, $Q'=L(R)$ for a smaller
rooted $R$; the expanded-diamond conclusion of
Lemma~\ref{lem:root-diamond} gives the contradiction.
These cases exhaust the root triangle. The final assertion follows
from bridge symmetrization and then minimum-order choice among rooted
obstructions. It is an existence reduction, not a classification of
all cubic graphs with bridges.
\end{proof}

\subsection{Certification through rooted order twenty-three}
We checked both normal forms at every edge of all 556,471 connected
simple cubic graphs of even order four through twenty, giving
33,079,650 indexed rooted tests. For the first form at
parent order twenty-two, the complete certificates from Section~25
supply the remaining cases. For a fragment of odd order $q$, each
new record retains either a root-out good set of size at least
$(q+1)/2$, or a root-out and a root-in good set whose sizes sum to
at least $q$. In either case the double satisfies its half bound.

The producer \texttt{root\_normal\_forms.cpp} and independent checker
\texttt{check\_root\_normal\_\allowbreak forms.cpp} reside in \texttt{exploratory/}.
The checker reconstructs each fragment and directly verifies its degrees
and selected-neighbour counts. The driver \texttt{verify\_round33.py}
checks generator and checker counts and the inherited certificate hashes.
Census completeness relies on nauty. Separately, exhaustive vertex-subset
enumeration on 84 small inputs verifies all four profile formulas,
including the exact single-expansion failure criterion.

\begin{corollary}[Finite bridge exclusion]\label{cor:root23-all}
Every connected rooted graph of the specified degree pattern on at most
23 vertices satisfies $d(Q)\geq|V(Q)|$. Any bridge join of two such
rooted graphs satisfies the original half bound. In particular:
\begin{enumerate}
\item every connected cubic counterexample with a bridge has at least
30 vertices;
\item a cubic graph on 46 vertices with a bridge splitting it into two
23-vertex components satisfies the half bound, with no restriction on
triangles at the bridge endpoints.
\end{enumerate}
\end{corollary}
\begin{proof}
If a rooted obstruction of order at most 23 existed, choose one of
minimum order. The normal-form theorem would give a parent of order
at most 22 in the first case or at most 20 in the second. All these
cases have checked positive certificates, a contradiction.
Mixed joins follow from bridge symmetrization.
Every rooted fragment has odd order at least five: orders one and
three cannot have the required degree-three vertices in a simple graph.
Thus deleting a bridge from a connected cubic graph of order at most
28 leaves two rooted components each of order at most 23. This proves
the first consequence; applying the same join statement to two
23-vertex components proves the second.
\end{proof}

\section{Restricting the distinguished edge and its parent profile}
The two forms in Theorem~\ref{thm:root-normal-form} can be restricted
further. The distinguished edge must lie on a cycle; this does not
require that the rest of the parent graph be bridgeless.

\begin{proposition}\label{prop:nonbridge-parent-edge}
In a minimum-order rooted obstruction in either normal form, the
replaced edge of the cubic parent is not a bridge.
\end{proposition}
\begin{proof}
Suppose $e=ab$ is a bridge of the parent $G$. Its deletion gives
rooted components $R_1,R_2$ of odd orders $q_1,q_2$. Write $A_i,B_i$
for their root-in and root-out optima and $M_i=\max(A_i,B_i)$.
These components are smaller than the putative rooted obstruction,
so minimality gives
\[
 d(R_i)=M_i+B_i\geq q_i.
\]
For the edge-subdivision form, the new root lies on no cycle. Its
root-out optimum is $M_1+M_2$. Since $2M_i\geq d(R_i)\geq q_i$
and $q_i$ is odd, $2M_i\geq q_i+1$. The double therefore has optimum
at least $q_1+q_2+2$, exceeding the required $q_1+q_2+1$.

For the rooted-triangle form, both joining edges from the new triangle
to $R_1,R_2$ are bridges. A root-in good set selects the triangle and
must avoid both attachment roots, giving
$A=3+B_1+B_2$. A root-out good set selects no vertex of that triangle,
giving $B=M_1+M_2$. Thus
\[
 A+B=3+d(R_1)+d(R_2)\geq3+q_1+q_2,
\]
exactly the required half bound for its double. Both cases contradict
rooted obstruction.
\end{proof}

\subsection{Exact tests on the cubic parent}
Let $G$ have order $2h$, let $e=uv$, and define
\[
 U=\cind(G-e),\qquad
 T=\max_{S\text{ good in }G,\ u,v\in S}|S|,\qquad
 W=\cind(G-\{u,v\}).
\]
As usual, an impossible constrained maximum is $-\infty$.
The subdivision form $Q_0$ has profile $(T+1,U,W)$.
The triangle form $Q_1=E(Q_0)$ consequently has profile
\[
 (A_1,B_1,C_1)=(W+3,\ \max(U,T+2),\ U).
\]
These are direct instances of the root expansion formula.

\begin{corollary}\label{cor:parent-triangle-obstruction}
If a minimum-order rooted obstruction has the triangle form, its
parent edge satisfies
\[
 U=h+1,\qquad T\leq h-1,\qquad W\leq h-2.
\]
Conversely, a cubic parent and edge satisfying these three conditions
would give a counterexample to the original conjecture by doubling
$Q_1$. Its order would be $4h+6$.
\end{corollary}
\begin{proof}
Minimality makes $Q_0$, which has order $2h+1$, non-obstructing.
Apply the exact single-expansion failure criterion of Section~26.1
to $(A,B,C)=(T+1,U,W)$ to get the three conditions.
Conversely, they give $B_1=h+1$ and $A_1\leq h+1$, so the double
has optimum at most $2h+2$, below its required $2h+3$.
\end{proof}

Every good set of the parent is either counted by $T$ or survives
the deletion of $e$. Consequently the displayed conditions imply
$\cind(G)\leq h+1$. If the parent satisfies the original half bound,
its optimum must therefore be $h$ or $h+1$. Every good set of size
at least $h$ must contain exactly one endpoint of $e$: containing
both contradicts the bound on $T$, and avoiding both contradicts
the bound on $W$.
If the parent is bipartite, every good set has even size, including
in $G-e$. The equality $U=h+1$ therefore forces $h$ odd, so a
bipartite parent of this minimum triangle form has order $2$ modulo $4$.
These necessary conditions do not establish the existence or
nonexistence of such a parent.

\subsection{Targeted search and independent checks}
Exhaustive vertex-subset checks on all 369 parent-edge cases from
connected simple cubic graphs through order ten verify both profile
identities directly. The bridge case in that census also verifies
the component formulas used in Proposition~\ref{prop:nonbridge-parent-edge}.

We then sought a rooted graph $Q_0$ of order 25 satisfying
$A\leq12$, $B\leq13$, $C\leq10$. Such a graph would yield an
original counterexample on 54 vertices after root expansion and
doubling. The search fixes the two root neighbours, forbids their
adjacency, and requires $Q_0-r$ connected so that the suppressed
parent edge is not a bridge. It permits deletion of at most six
edges of a fixed seed, with degree constraints forcing the same
number of replacements. The seed comes from subdividing a non-bridge
edge of a retained sharp 24-vertex cubic parent.

A counterexample-guided satisfiability search proposes graphs and
exhaustively enumerates good sets of size at least eleven. Every
violating set adds a constraint excluding that selected set from
being good in subsequent candidates. The three thresholds detect
all violations of the stated target bounds. Both the initial
rooted double and its expanded double are also tested, to avoid
missing an earlier counterexample. A separate cycle-space verifier
checks the exact profiles of all retained connected candidates.
The search is restricted to a labelled neighbourhood of one seed;
a timeout or an unknown solver result does not exhaust this domain.
The source is \texttt{exploratory/root\_triangle\_synthesis.py},
and the independent verifier is \texttt{verify\_round34.py}.

The run reached its wall-time stopping condition after 407.4 seconds,
with 12 retained candidates and 1,919 exclusion constraints, including
160 initial constraints. The wall-time condition is checked between
iterations. Independent enumeration of all 8,192 cycle-space vectors
for each candidate confirms its exact profile. The minimum doubled
optima among these candidates are 30 before expansion and 32 after
expansion, compared with required half bounds of 25 and 27.
No candidate is a counterexample. The remaining search region was
not exhausted, and no general sufficiency proof follows from the run.

\clearpage
\begin{thebibliography}{99}
\bibitem{Speyer}
D.~E.~Speyer, \emph{Cyclically Orientable Graphs},
preprint, 2005.
\href{https://arxiv.org/abs/math/0511233}{arXiv:math/0511233}.
\bibitem{Gutin}
G.~Gutin, \emph{Note on Perfect Forests}, Journal of Graph Theory
\textbf{82}, no.~3 (2016), 233--235.
\href{https://doi.org/10.1002/jgt.21897}{doi:10.1002/jgt.21897}.
\bibitem{Edmonds}
J.~Edmonds, \emph{Maximum matching and a polyhedron with
$0,1$-vertices}, Journal of Research of the National Bureau of
Standards, Section B \textbf{69B} (1965), 125--130.
\href{https://doi.org/10.6028/jres.069B.013}{doi:10.6028/jres.069B.013}.
\bibitem{Petersen}
J.~Petersen, \emph{Die Theorie der regul\"aren graphs},
Acta Mathematica \textbf{15} (1891), 193--220.
\href{https://doi.org/10.1007/BF02392606}{doi:10.1007/BF02392606}.
\bibitem{HJLS}
M.~A.~Henning, F.~Joos, C.~L\"owenstein and T.~Sasse,
\emph{Induced Cycles in Graphs},
Graphs and Combinatorics \textbf{32} (2016), 2425--2441.
\href{https://arxiv.org/abs/1406.0606}{arXiv:1406.0606}.
\bibitem{KSS}
M.~Knor, J.~Sedlar and R.~\v{S}krekovski,
\emph{Counterexamples to two conjectures on $(1,2)$-domination
in cubic graphs}, 2026.
\href{https://arxiv.org/abs/2608.17851}{arXiv:2608.17851}.
\bibitem{RRT}
G.~Radi\'c, G.~Radi\'c and A.~Tepeh,
\emph{Structural aspects of $1$-limited domination in cubic graphs},
2026.
\href{https://arxiv.org/abs/2610.26259}{arXiv:2610.26259}.
\bibitem{Sage}
The Sage Developers, \emph{SageMath graph generators: small graphs},
source module \texttt{smallgraphs.py},
accessed October 27, 2026.
\href{https://doc.sagemath.org/html/en/reference/graphs/sage/graphs/generators/smallgraphs.html}{SageMath reference documentation}.
\bibitem{GJMMMU}
J.~Goedgebeur, J.~Jooken, E.~M\'a\v{c}ajov\'a, D.~Mattiolo,
G.~Mazzuoccolo and N.~Ulyanov,
\emph{Disproving the Petersen Coloring Conjecture: Theoretical Analysis
and an Infinite Family of Counterexamples}, 2026.
\href{https://arxiv.org/abs/2608.10028v4}{arXiv:2608.10028v4}.
\bibitem{HoGCubic}
House of Graphs, \emph{Connected cubic graphs}, census table,
accessed October 28, 2026.
\href{https://houseofgraphs.org/Cubic}{houseofgraphs.org/Cubic}.
\end{thebibliography}
\end{document}
